%
% architecture.tex
%
% Quarry Technical Architecture Specification
% Written for principal engineers evaluating the system.
%
% Compile: pdflatex architecture.tex
%

\documentclass[a4paper,10pt]{article}
\usepackage[margin=1in]{geometry}
\usepackage{amsmath}
\usepackage{booktabs}
\usepackage{listings}
\usepackage{xcolor}
\usepackage{enumitem}
\usepackage{longtable}
\usepackage[colorlinks=true,linkcolor=blue,citecolor=blue,urlcolor=blue]{hyperref}

\lstset{
  basicstyle=\ttfamily\small,
  keywordstyle=\color{blue!80!black},
  commentstyle=\color{green!50!black},
  stringstyle=\color{red!60!black},
  breaklines=true,
  frame=single,
  framerule=0.3pt,
  rulecolor=\color{gray!50},
  backgroundcolor=\color{gray!5},
  columns=fullflexible,
}

\newcommand{\modpath}[1]{\texttt{src/quarry/#1}}
\newcommand{\tool}[1]{\texttt{#1}}
\newcommand{\field}[1]{\texttt{#1}}

\begin{document}

\title{Quarry: Technical Architecture Specification}
\author{Local Semantic Search for AI Agents}
\date{September 2026}
\hypersetup{
  pdftitle={Quarry: Technical Architecture Specification},
  pdfauthor={Punt Labs},
  pdfsubject={Architecture},
  pdfcreator={pdflatex}
}
\maketitle

\tableofcontents
\newpage

% ===================================================================
\section{Overview}
% ===================================================================

Quarry is a local semantic search engine.  It ingests documents in
20+ formats, embeds them with a local ONNX model, stores vectors in
LanceDB, and serves semantic search to AI agents (Claude Code, Claude
Desktop) and a macOS menu bar app.

The system is \textbf{daemon-first} (DES-031 v2.2, ACCEPTED and
complete): a single \texttt{quarryd} process holds the engine, and the
CLI, MCP, and library are pure clients that reach it over a versioned
wire protocol --- the daemon is the intermediary service layer.  The
engine runs only in \texttt{quarryd} (the \texttt{quarry serve}
subcommand is retired); loopback requests are authenticated with a
mode-0600 \texttt{serve.token} resolved through \texttt{ClientConfig};
and every client surface --- the CLI, the \texttt{quarry mcp} stdio
server, hooks, and the library --- reaches the daemon through one
\texttt{QuarryClient} over the \texttt{/v1} REST contract.  No client
process imports or constructs the engine; an import-linter contract and
a runtime engine-sabotage test enforce that boundary structurally in CI
(the epic \texttt{quarry-ma6f}/\texttt{quarry-ynvs} is closed).  The
always-on watch/index loop that daemon-first exists to enable (DES-045)
keeps every registered directory's index fresh with no human action.

\subsection{Design Principles}

\begin{enumerate}[label=\textbf{P\arabic*}.]
  \item \textbf{Daemon-first.}  The \texttt{quarryd} process is
    the single engine (embedding model, LanceDB, pipeline, sync
    registry); the CLI, MCP, and library are clients over a
    versioned REST contract (DES-031 v2.2, ACCEPTED and complete).
    The prior \emph{library-first} principle --- the core library did
    all work and the surfaces were thin wrappers --- described the
    as-built state DES-031 v2 replaced.
  \item \textbf{Local by default.}  Everything runs on the user's
    machine.  No API keys, no cloud accounts.  The embedding model
    downloads once ($\sim$120\,MB, int8 quantized).
  \item \textbf{One daemon, many clients.}  A single
    \texttt{quarryd} process loads the embedding model once and
    serves all sessions over the REST wire; supervised by launchd
    \texttt{KeepAlive} / systemd \texttt{Restart=always}.
  \item \textbf{Fire-and-forget I/O.}  Side-effect MCP tools return
    immediately and process in a bounded thread pool.  Search is
    synchronous.
  \item \textbf{Named databases.}  Each database is fully isolated:
    its own LanceDB directory, sync registry, and vector index.
    Work/personal separation without configuration files.
\end{enumerate}

% ===================================================================
\section{System Architecture}
% ===================================================================

\subsection{Deployment Topology}

\textbf{Single machine (localhost TLS):}
\begin{lstlisting}
                       stdio               HTTPS /v1 (TLS)
Claude Code <------------------> quarry mcp <-------------------> quarryd
             MCP JSON-RPC        (QuarryClient) pinned CA cert    (one daemon)
\end{lstlisting}

\textbf{Remote server (GPU):}
\begin{lstlisting}
Mac (client)                                   Linux (server)
                       stdio               HTTPS /v1 (TLS)
Claude Code <------------------> quarry mcp <---------------------> quarryd
             MCP JSON-RPC        (QuarryClient) pinned CA cert     (GPU daemon)

CLI -----> HTTPS GET/POST /v1 -----> quarryd
           Bearer + pinned CA
\end{lstlisting}

Quarry's own Claude Code plugin runs \texttt{quarry mcp} directly (a
FastMCP stdio server whose tools are \texttt{QuarryClient} calls over the
daemon's \texttt{/v1} REST API); it no longer routes through
\texttt{mcp-proxy} (DES-031 v2.2).  Because a single \texttt{quarryd}
holds the resident engine, every Claude Code tab shares one warm model
rather than spawning a Python process that loads $\sim$200\,MB of its
own.  Installed services use TLS with TOFU certificate pinning (DES-020)
--- even on localhost.  (\texttt{mcp-proxy} is an MCP stdio$\leftrightarrow$daemon
transport, not a REST-API client; \texttt{quarry install} still fetches it
because it remains a supported transport for other MCP clients, but it is
no longer on quarry's own MCP path --- the plugin runs \texttt{quarry mcp}
directly.)

When a remote server is configured (DES-021), the CLI routes
all data commands over HTTPS to the server, with only authentication
and administration commands remaining local.  \texttt{find},
\texttt{show}, \texttt{status}, \texttt{ingest}, \texttt{remember},
\texttt{delete}, \texttt{sync}, all \texttt{list} subcommands, and
directory \texttt{register}/\texttt{deregister} route remotely.
\texttt{use} is local by design: it sets the client's persistent default
database, but the remote server is fixed to the database it was started
with, so switching databases is a client-side setting rather than a
remote call.

\subsection{Surface Summary}

\begin{center}
\begin{tabular}{llp{7cm}}
\toprule
\textbf{Surface} & \textbf{Transport} & \textbf{Purpose} \\
\midrule
CLI         & \texttt{QuarryClient} over HTTPS \texttt{/v1} & Interactive use from terminal \\
MCP server  & stdio (\texttt{quarry mcp}, a \texttt{QuarryClient}) & Claude Code/Desktop tools \\
HTTP server & REST JSON \texttt{/v1} (FastAPI in \texttt{quarryd}) & CLI, MCP, menu bar app \\
Plugin      & Hook shell commands (\texttt{quarry-hook}) & Auto-register; capture web fetch/search,
              opt-in reads, compaction, session end, subagent stop; memory distillation \\
\bottomrule
\end{tabular}
\end{center}

% ===================================================================
\section{Module Responsibilities}
% ===================================================================

\subsection{Data Layer (\modpath{db/})}

The database facade composes five subsystems behind a single entry
point (\texttt{Database}), so callers pass one object rather than a raw
LanceDB connection.

\begin{center}
\begin{tabular}{llp{7cm}}
\toprule
\textbf{Module} & \textbf{File} & \textbf{Responsibility} \\
\midrule
Facade        & \modpath{db/facade.py}        & \texttt{Database} ---
                single entry point composing the store, search, catalog,
                schema, and optimizer subsystems (DES-031). \\
Chunk store   & \modpath{db/chunk\_store.py}   & Insert,
                progressive-flush, delete, and count operations on the
                \texttt{chunks} table. \\
Chunk search  & \modpath{db/chunk\_search.py}  & Vector-similarity
                search under the cosine metric. \\
Chunk catalog & \modpath{db/chunk\_catalog.py} & List documents, list
                collections, and fetch page text. \\
Optimizer     & \modpath{db/optimizer.py}      & Table compaction, FTS
                index rebuild, and per-collection indexing. \\
Schema        & \modpath{db/schema.py}         & Chunks-table schema and
                idempotent migration (agent-memory columns, FTS index). \\
Storage       & \modpath{db/storage.py}        & LanceDB connection
                management (discovery and size reporting removed --- DES-049). \\
Connection    & \modpath{db/connection.py}     & Self-recycling LanceDB
                connection that bounds file-descriptor growth. \\
Chunk table   & \modpath{db/chunk\_table.py}   & Chunks-table lifecycle:
                open, migrate, create-on-demand. \\
\bottomrule
\end{tabular}
\end{center}

Supporting types: \modpath{models.py} (immutable \texttt{PageContent},
\texttt{Chunk}, \texttt{PageAnalysis} dataclasses), \modpath{types.py}
(structural \texttt{Protocol}s for LanceDB, OCR, and embedding
backends), \modpath{results.py} (\texttt{SearchResult},
\texttt{SearchFilter}, \texttt{WORST\_CASE\_DISTANCE}), and
\modpath{config.py} (\texttt{Settings}, \texttt{resolve\_db\_paths()},
ONNX model constants).

\subsection{Ingestion Pipeline (\modpath{ingestion/})}

\begin{center}
\begin{tabular}{llp{7cm}}
\toprule
\textbf{Module} & \textbf{File} & \textbf{Responsibility} \\
\midrule
Pipeline    & \modpath{ingestion/pipeline.py}    & File orchestrator:
              resolve the format strategy, extract, then chunk/embed/store
              through the funnel.  Entry points \texttt{ingest\_document}
              and \texttt{plan\_file\_chunks}; re-exports
              \texttt{SUPPORTED\_EXTENSIONS}. \\
Format strategies & \modpath{ingestion/format\_strategies.py} &
              \texttt{resolve\_strategy(suffix)} returns a
              \texttt{FormatStrategy} (PDF, image, or a
              \texttt{TextLikeFormat} row from
              \modpath{ingestion/text\_format.py}); owns
              \texttt{SUPPORTED\_EXTENSIONS}. \\
Funnel      & \modpath{ingestion/chunk\_store\_funnel.py} & The single
              convergence point for chunking, embedding, and storing pages
              (the DES-036 choke point every ingest path shares). \\
Chunker     & \modpath{ingestion/chunker.py}     & Split page text into
              overlapping chunks; respects page and heading boundaries. \\
Web ingest  & \modpath{ingestion/web\_ingest.py} & \texttt{WebIngest}:
              inline content (\texttt{ingest\_content}) and URL
              (\texttt{ingest\_url}) ingest --- extract, scrub when the
              request carries a \texttt{content\_scrubber}, funnel through
              the store. \\
Bulk / sitemap & \modpath{ingestion/bulk\_ingest.py},
              \modpath{ingestion/sitemap\_ingest.py} & Threaded bulk fan-out
              and sitemap-driven URL ingestion. \\
Web fetch   & \modpath{ingestion/web\_fetch.py}  & Fetch a URL over
              HTTP(S) and return validated HTML text. \\
Ingest stats & \modpath{ingestion/ingest\_stats.py} & Per-format ingest
              counters merged into an ingest result. \\
\bottomrule
\end{tabular}
\end{center}

Format extraction lives in \modpath{extractors/}; embedding and OCR
provider selection in \modpath{ingestion/provider.py},
\modpath{ingestion/backends.py}, and \modpath{ingestion/ocr\_local.py}.

\subsection{Retrieval Seam (\modpath{retrieval/})}

Hybrid search runs through one parameterizable seam shared by every
surface and by the eval harness, so the production path and the
committed baseline cannot drift.

\begin{center}
\begin{tabular}{llp{7cm}}
\toprule
\textbf{Module} & \textbf{File} & \textbf{Responsibility} \\
\midrule
Service    & \modpath{retrieval/service.py}   & \texttt{SearchService}
             --- the single call every surface (CLI, HTTP, MCP) makes
             for search.  Kills remote/local divergence (bug class 3). \\
Hybrid     & \modpath{retrieval/hybrid.py}    & \texttt{HybridRetriever}
             --- vector + BM25 channels fused with RRF, behind a
             config. \\
Config     & \modpath{retrieval/config.py}    & \texttt{RetrievalConfig}
             --- frozen retrieval knobs (RRF $k$, over-fetch, cosine
             metric, decay rate, reranker).  Defaults reproduce shipped
             search bit-for-bit. \\
Protocols  & \modpath{retrieval/protocols.py} & \texttt{Retriever} and
             \texttt{Reranker} structural interfaces (one method
             each). \\
Fusion     & \modpath{retrieval/fusion.py}    & \texttt{RrfFusion} ---
             Reciprocal Rank Fusion with optional temporal decay. \\
Reranker   & \modpath{retrieval/reranker.py}  & \texttt{NullReranker}
             --- pass-through Null Object when reranking is off. \\
\bottomrule
\end{tabular}
\end{center}

\subsection{Format Processors (\modpath{extractors/})}

Each extractor implements a common \texttt{Protocol}
(\modpath{extractors/protocol.py}) and converts a source format into
\texttt{list[PageContent]}:

\begin{center}
\begin{tabular}{llp{6cm}}
\toprule
\textbf{Module} & \textbf{File} & \textbf{Formats} \\
\midrule
PDF           & \modpath{extractors/pdf\_extractor.py}          & PDF
                (text via PyMuPDF, OCR fallback for image pages) \\
Image         & \modpath{extractors/image\_extractor.py}        & PNG,
                JPG, TIFF, BMP, WebP (OCR) \\
Text          & \modpath{extractors/text\_extractor.py}         & TXT,
                MD, LaTeX, DOCX \\
Code          & \modpath{extractors/code\_extractor.py}         & 30+
                languages (tree-sitter section splitting) \\
HTML          & \modpath{extractors/html\_extractor.py}         & HTML
                (boilerplate stripping, Markdown conversion) \\
Spreadsheet   & \modpath{extractors/spreadsheet\_extractor.py}  & XLSX,
                CSV (LaTeX tabular serialization) \\
Presentation  & \modpath{extractors/presentation\_extractor.py} & PPTX
                (slide-per-chunk with notes) \\
\bottomrule
\end{tabular}
\end{center}

\subsection{Backends}

\begin{center}
\begin{tabular}{llp{7cm}}
\toprule
\textbf{Module} & \textbf{File} & \textbf{Responsibility} \\
\midrule
Factory     & \modpath{ingestion/backends.py}  & Thread-safe factory
              with double-checked locking.  \texttt{get\_ocr\_backend()}
              and \texttt{get\_embedding\_backend()} return cached
              singletons. \\
OCR         & \modpath{ingestion/ocr\_local.py} & Local OCR via RapidOCR
              (offline, no setup). \\
Embeddings  & \modpath{embeddings.py}          & Local ONNX embedding
              (snowflake-arctic-embed-m-v1.5, 768-dim, 512 tokens).
              Vectors are L2-normalized to unit length. \\
Provider    & \modpath{ingestion/provider.py}  & ONNX Runtime execution
              provider auto-detection (CUDA/FP16 or CPU/int8). \\
\bottomrule
\end{tabular}
\end{center}

\subsection{Capture and Redaction}

Every automatic capture --- session transcripts (PreCompact, SessionEnd,
SubagentStop), WebFetch pages, WebSearch digests, and opt-in Read
captures --- is scrubbed at write time before any bytes reach a
git-tracked (and pushable) surface (DES-036), and again server-side by
the daemon's \texttt{ScrubbedIngestJob} before a chunk is stored
(DES-041).

\begin{center}
\begin{tabular}{llp{7cm}}
\toprule
\textbf{Module} & \textbf{File} & \textbf{Responsibility} \\
\midrule
Capture writer & \modpath{capture.py}       & \texttt{CaptureWriter} ---
                 the single choke point every local \texttt{.md} capture
                 producer routes through.  Scrubs, then atomically writes. \\
Scrubber       & \modpath{scrub.py}         & \texttt{Scrubber} ---
                 ordered redaction passes (secrets, paths, emails,
                 hostname, profanity).  Idempotent. \\
Scrub rules    & \modpath{scrub\_rules.py}  & Secret-detection regex
                 catalog applied by the scrubber. \\
Capture URL    & \modpath{capture\_url.py}  & \texttt{CaptureUrl} ---
                 strips userinfo/query/fragment from a captured URL's
                 metadata form. \\
Web capture    & \modpath{web\_capture.py}  & Parse a PostToolUse
                 WebFetch payload into URL + fetched content. \\
Transcript     & \modpath{transcript.py}    & The session transcript
                 file and its derived identifiers. \\
Session capture & \modpath{session\_transcript.py} &
                 \texttt{SessionTranscriptCapture} --- archive the raw
                 JSONL under \texttt{\~{}/.punt-labs/quarry/sessions/},
                 write the scrubbed \texttt{.md}, POST the
                 \texttt{CaptureIngestRequest}; shared by PreCompact,
                 SessionEnd, and SubagentStop. \\
Subagent capture & \modpath{subagent\_capture.py},
                 \modpath{subagent\_report.py} & \texttt{SubagentCapture}
                 --- the raw transcript row plus the subagent's final
                 report filed as an \texttt{observation} in
                 \texttt{memory-<handle>} (DES-055). \\
Search / read  & \modpath{web\_search\_capture.py},
                 \modpath{read\_capture.py} & Parse a WebSearch result
                 digest; the four-check admission filter for the opt-in
                 Read capture. \\
Daemon sender  & \modpath{daemon\_capture.py} & \texttt{DaemonCaptureSender}
                 --- the hooks' one bounded door to \texttt{/v1/capture},
                 \texttt{/v1/ingest}, and \texttt{/v1/remember}. \\
Backfill       & \modpath{backfill.py}      & Backfill historical
                 session transcripts through the capture writer. \\
\bottomrule
\end{tabular}
\end{center}

\subsection{Fetch Safety (SSRF)}

Every outbound fetch --- the leaf-page fetch and the sitemap crawl ---
runs the address-safety gate (\texttt{UrlSafetyCheck}) against the
\emph{resolved} address on \emph{every} hop, fail-closed, before the
socket opens (DES-043).  The route boundary validates only the initial
caller-supplied URL; the gap this closes is that redirect targets and
sitemap-discovered URLs never reached that gate --- an authenticated
caller could otherwise reach \texttt{169.254.169.254}, loopback, or a
private range via a redirect from a public server or an internal URL
listed in a sitemap.

\begin{center}
\begin{tabular}{llp{7cm}}
\toprule
\textbf{Module} & \textbf{File} & \textbf{Responsibility} \\
\midrule
URL safety     & \modpath{url\_safety.py}   & \texttt{UrlSafetyCheck}
                 --- classifies a URL's resolved address (rejects
                 link-local, loopback, RFC-1918, CGNAT, unspecified,
                 IPv4-mapped-IPv6, and the NAT64 well-known prefix;
                 checks all \texttt{getaddrinfo} records).  In core, so
                 the daemon route and the fetch layer share one
                 classifier. \\
Redirect guard & \modpath{ingestion/ssrf\_redirect.py} &
                 \texttt{SsrfGuardedRedirectHandler} plus the shared
                 \texttt{GUARDED\_OPENER}: re-gates every 30x
                 \texttt{Location} before it is opened. \\
Sitemap client & \modpath{sitemap\_web\_client.py} &
                 \texttt{GatedSitemapWebClient} --- USP's
                 \texttt{AbstractWebClient}, gating every sitemap-index,
                 \texttt{robots.txt}, and sub-sitemap fetch at every
                 recursion depth (a block is non-retryable, so USP
                 never connects). \\
Pinned connection & \modpath{ingestion/pinned\_connection.py} &
                 \texttt{PinnedHTTP(S)Connection} --- performs one
                 \texttt{getaddrinfo}, validates every record, and connects
                 the socket to a validated IP literal from that same set; no
                 second resolution exists, so the gate-vs-connect DNS-rebind
                 window is eliminated by construction (DES-044).  Installed
                 into the shared \texttt{GUARDED\_OPENER}, so every fetch path
                 and redirect hop inherits the pin. \\
\bottomrule
\end{tabular}
\end{center}

A single \texttt{getaddrinfo} per connection eliminates the
gate-vs-socket DNS-rebind TOCTOU by construction: the socket connects
to an address the gate has just validated, never an independent
re-resolution (DES-044, closing quarry-kmzo/quarry-ljym).  The pin
narrows the \emph{address}, not the \emph{trust} --- \texttt{self.host}
is never mutated, so SNI, the \texttt{Host} header, and certificate
verification stay on the hostname (never the pinned IP), and the
public-fetch context keeps the system trust store, kept separate from
the daemon-RPC pinned-CA context.

\subsection{Sync and Registry}

\begin{center}
\begin{tabular}{llp{7cm}}
\toprule
\textbf{Module} & \textbf{File} & \textbf{Responsibility} \\
\midrule
Sync       & \modpath{sync.py}           & Directory sync: discover
             files, compute delta (new/changed/deleted),
             ingest/delete accordingly. \\
Discovery  & \modpath{sync\_discovery.py} & Walk, filter, and hash files
             for the sync delta. \\
Ingest     & \modpath{sync\_ingest.py}    & Progressive per-collection
             ingest (producer/consumer with within-file resume). \\
Registry   & \modpath{sync\_registry.py}  & SQLite registry of watched
             directories and per-file state. \\
Planner    & \modpath{sync\_planner.py}   & Compute the disk-vs-registry
             plan for one collection (TOCTOU-resilient; the one
             re-embed rule lives in \modpath{sync\_change.py}). \\
Lock       & \modpath{sync\_lock.py}      & \texttt{SyncLock} --- the
             kernel-\texttt{flock} single-flight guard around the detached
             background sync launch. \\
Finalize   & \modpath{sync\_finalize.py}  & Post-sync index rebuild,
             optimize, capture-shadow push, and GC. \\
Resume     & \modpath{sync\_resume.py}    & Within-file resume policy:
             watermark gate and partial-hash marks (DES-034). \\
\bottomrule
\end{tabular}
\end{center}

\subsection{Always-on Watch/Index Loop (DES-045)}

Continuous incremental indexing --- the first-class daemon-first rationale
(DES-031) --- is a daemon-owned \modpath{daemon/watch\_loop.py}
\texttt{WatchLoop} that keeps every registered directory's index fresh with no
human action. It is a \emph{producer only}: it watches the filesystem,
debounces edit bursts, and submits index work to the existing DES-042 queue,
writing no LanceDB table directly. So it inherits the whole serialization
stack unchanged --- DES-042 per-table FIFO $\to$ DES-034
\texttt{ProgressiveIndexer}/\texttt{\_write\_lock} $\to$ DES-026 WAL --- and the
queue's routing key becomes \texttt{(database, collection)}, extending
single-writer-per-table across the roster. Three producers (initial scan on
start, the live watch, an explicit \texttt{quarry sync}) all feed the one queue
and thus cannot race.

\begin{center}
\begin{tabular}{llp{6.6cm}}
\toprule
\textbf{Module} & \textbf{File} & \textbf{Responsibility} \\
\midrule
Watch loop  & \modpath{daemon/watch\_loop.py}   & The \texttt{WatchLoop}:
              per-database watch sets + observer lifecycle, initial scan,
              start/stop watching, \texttt{request\_scan}.  Producer into
              the queue. \\
Reconcile   & \modpath{daemon/watch\_reconcile.py} & The
              \texttt{watch\_safety\_scan\_s} roster reconcile: re-enumerate,
              start watching new databases, re-submit shed scans, run the
              orphan sweep (DES-045a). \\
Submit      & \modpath{daemon/watch\_submit.py} & Turns a debounced delta
              into per-file or bulk \texttt{IngestUnit}s and enqueues them. \\
Roster      & \modpath{daemon/watch\_roster.py} & Per-database connections +
              watch sets enumerated from the on-disk roster (trust boundary:
              only roster DBs are opened). \\
Scratch guard & \modpath{scratch\_paths.py}    & \texttt{ScratchGuard}: refuses
              a temp/scratch/VCS root and skips gitignored paths at every
              producer seam (DES-045b). \\
Removal     & \modpath{daemon/registration\_lifecycle.py} & Deregister /
              subsume / retain: marks a collection in
              \texttt{pending\_purge\_collections} on non-keep-data removal;
              records retained-collection identity (DES-045a). \\
Purge       & \modpath{daemon/purge\_service.py} & The orphan sweep +
              \texttt{CollectionPurgeJob}: deletes only marked collections
              still holding chunks, re-checked under its own transaction. \\
Route key   & \modpath{daemon/route\_key.py}   & \texttt{RouteKey =
              (database, collection)} --- the queue routing key that extends
              single-writer-per-table across the roster. \\
Events      & \modpath{daemon/fs\_events.py}    & \texttt{FsEvent} +
              \texttt{FsEventSource} protocol seam (tests inject synthetic
              events). \\
Watchdog    & \modpath{daemon/fs\_watchdog.py}  & The sole \texttt{watchdog}
              importer: FSEvents (macOS) / inotify (Linux) with a
              \texttt{PollingObserver} fallback (on \texttt{ENOSPC}). \\
Debounce    & \modpath{daemon/debounce.py}      & Per-path debounce window +
              max-delay cap + bulk-threshold coalescing. \\
Index jobs  & \modpath{daemon/index\_jobs.py}   & \texttt{FileIndexJob} /
              \texttt{DocumentDeleteJob} / \texttt{CollectionSyncJob}
              (\texttt{IngestUnit}s). \\
Finalize    & \modpath{daemon/finalize\_job.py} & \texttt{CollectionFinalizeJob}
              (\texttt{SyncFinalizer}): the lone post-quiescence
              \texttt{optimize()} + FTS rebuild, coalesced by
              \texttt{FinalizeThrottle} (DES-045c). \\
\bottomrule
\end{tabular}
\end{center}

The loop feeds the same DES-042 serialized queue (\modpath{daemon/ingest\_queue.py},
\modpath{daemon/collection\_worker.py}): one resident FIFO worker per
\texttt{RouteKey} (one in-flight \texttt{progressive\_insert} per LanceDB
table) behind a global embed \texttt{Semaphore} hard-clamped to~1 (DES-032:
more than one buys no matmul parallelism behind the shared ONNX session mutex
and re-adds arena contention).  Admission is non-blocking and bounded --- a
full queue returns 503 and the event re-arms with capped backoff, never a
silent drop.

A small delta re-indexes as per-file jobs (embed gate released between files
for cross-collection fairness); a burst above \texttt{watch\_bulk\_threshold}
(default~50) collapses to one \texttt{CollectionSyncJob} (DES-034's size-gated
flush keeps the fragment budget).  Per-file jobs \emph{never} rebuild the FTS
index --- a lone finalize job runs post-quiescence per \texttt{RouteKey}, so
the daemon's persistent connection does not reopen the quarry-0dss descriptor
leak (\texttt{create\_fts\_index(replace=True)} pins deleted-file readers the
LanceDB Rust core never evicts); a resource-invariant test proves the fd count
plateaus across multiple databases.

\textbf{Removal lifecycle and orphan sweep (DES-045a).}  Deregister,
subsumption, and orphan cleanup were hardened under a Z model
(\texttt{docs/spec/watch\_lifecycle.tex}, fuzz-clean and ProB-checked against
ten invariants I1--I10).  A collection is marked in a durable
\texttt{pending\_purge\_collections} table --- in the \emph{same} transaction as
the \texttt{directories} delete --- only on a non-keep-data deregister or a
subsume eviction; the periodic reconcile's orphan sweep then purges only
\emph{marked} collections that still hold chunks, re-checking under its own
transaction at execution time.  Non-directory collections (captures,
\texttt{memory-*}, the \texttt{default} \texttt{remember} target, URL-host
collections) are never marked, so the sweep is \emph{structurally unable} to
delete them (I9) --- the invariant whose earlier absence, in a model that
wrongly assumed every collection was a directory registration, would have wiped
all captures and agent memories on the 5-minute default timer.  A
\texttt{--keep-data} disable archives the collection's identity
(\texttt{retained\_collections} with \texttt{original\_directory}); only the
\emph{same} directory may re-adopt it (I7), and re-adoption triggers a reconcile
that prunes chunks for files deleted while it was disabled.  A different
directory always gets a fresh, distinct collection, and the fresh-name picker
avoids every chunk-bearing name so the auto path is merge-proof (I10).

\textbf{Watch scope (DES-045b).}  The \texttt{ScratchGuard} predicate, consulted
at every producer seam (initial scan, live watch, reconcile, explicit
\texttt{quarry sync}), refuses a root that is (a)~an OS temp root
(\texttt{/tmp}, \texttt{/private/tmp}, \texttt{/var/tmp},
\texttt{/private/var/tmp}, \texttt{/var/folders}, compared casefolded so a case
variant on APFS cannot slip past) or (b)~a repo's own scratch, anchored on
\texttt{<gitroot>/.tmp} for \emph{any} ancestor git repo; below a permitted
root it skips gitignored paths and an always-skip set (\texttt{.git},
\texttt{.beads}, \texttt{.venv}, \texttt{node\_modules}, \texttt{\_\_pycache\_\_},
\texttt{dist}, \texttt{build}, caches).  Resolution is fail-closed \emph{per
root} --- a \texttt{resolve()} that raises skips that one registration and logs,
never aborting the scan.  The guard fires at \emph{watch} time, not only at
register time, so an already-registered temp root is refused on daemon restart
with no registry edit.  This closes the hole where a stale \texttt{/private/tmp}
registration made the loop re-index the churning system temp dir and never reach
quiescence (283 CPU-minutes observed before the guard shipped).

\textbf{CPU/compaction envelope (DES-045c).}  DES-045's ``bounded by
construction'' bounded fds, memory, and the write-serializer, but LanceDB sizes
its own tokio \emph{compute} runtime to \texttt{num\_cpus}, and a single finalize
(\texttt{optimize()} + native FTS rebuild) on that uncapped pool seized every
core (291--403\,\% CPU observed).  Two structural mitigations bound it: (1)~\
\texttt{ThreadConfig.apply\_env\_limits()} (\modpath{thread\_config.py}) clamps
\texttt{LANCE\_CPU\_THREADS}/\texttt{LANCE\_IO\_THREADS} (cap~2) before the first
\texttt{lancedb.connect} in \texttt{DaemonServer.run()}, via a fail-closed clamp
with a per-variable floor (a stale \texttt{LANCE\_CPU\_THREADS=32} is forced down
to the cap, and \texttt{=1} --- which stalls lance's runtime --- is raised to the
floor of~2, so floor equals ceiling for the LANCE pools); (2)~\texttt{FinalizeThrottle}
coalesces the finalize to at most one per \texttt{watch\_optimize\_min\_interval\_s}
(default~30\,s) \emph{per database}, with a trailing timer so FTS still catches
up after churn settles.  The vector channel is never stale (files index
immediately); the reconcile is the unthrottled backstop.

\subsection{Surfaces}

\begin{center}
\begin{tabular}{llp{7cm}}
\toprule
\textbf{Module} & \textbf{File} & \textbf{Responsibility} \\
\midrule
CLI         & \modpath{\_\_main\_\_.py} + \modpath{cli\_*.py} & Typer
              CLI with rich formatting.  A thin \texttt{QuarryClient}
              client: build request $\to$ call daemon $\to$ format
              output.  Command groups split into \modpath{cli\_search},
              \modpath{cli\_ingest}, \modpath{cli\_documents},
              \modpath{cli\_sync}, \modpath{cli\_project},
              \modpath{cli\_remote}, \modpath{cli\_maintenance},
              \modpath{cli\_captures}, \modpath{cli\_missions}; shared
              output plumbing in \modpath{cli\_formatters}. \\
MCP Server  & \modpath{mcp\_server.py}, \modpath{mcp\_missions.py} &
              \texttt{quarry mcp}: a FastMCP \emph{stdio} server whose
              tool bodies (the search/ingest/management tools on
              \texttt{McpTools}, plus \texttt{missions\_sync} on the sibling
              \texttt{MissionTools}, all behind \modpath{mcp\_guard.py}'s
              \texttt{ToolGuard}) are
              thin \texttt{QuarryClient} calls over the daemon's
              \texttt{/v1} API.  Imports \emph{no} engine --- loads zero
              LanceDB/ONNX (DES-031 v2.2). \\
Client tier & \modpath{client/}          & \texttt{QuarryClient} +
              injectable REST transport, typed \texttt{QuarryError}
              hierarchy, \texttt{ClientConfig} target/token resolution.
              Replaces the deleted \texttt{RemoteClient}. \\
Wire API    & \modpath{api/}             & The \texttt{/v1} contract:
              one Pydantic request/response model per operation +
              a uniform \texttt{ErrorBody}.  Drift is an import-time
              type error (the structural kill of bug class 3). \\
Hooks       & \modpath{hooks.py}, \modpath{hooks\_agent.py} & Claude
              Code plugin hooks: session-start (auto-register),
              post-web-fetch, pre-compact (\modpath{hooks.py});
              session-end, post-web-search, post-read, subagent-stop
              (\modpath{hooks\_agent.py}).  Pure clients
              (\texttt{QuarryClient}); all fail-open. \\
Hook Entry  & \modpath{\_hook\_entry.py} & Lightweight hook dispatcher
              that bypasses full CLI import chain.  Dispatches
              via \texttt{sys.argv} with lazy imports
              ($\sim$0.1\,s vs $\sim$1.5\,s for full CLI). \\
Remote      & \modpath{remote.py}       & Remote login config:
              TOFU CA cert fetch and storage, pinned-CA connection
              validation. \\
Daemon      & \modpath{service.py}      & \texttt{quarryd} lifecycle
              via \texttt{launchd}/\texttt{systemd}. \\
\bottomrule
\end{tabular}
\end{center}

\subsection{Support}

\begin{center}
\begin{tabular}{llp{7cm}}
\toprule
\textbf{Module} & \textbf{File} & \textbf{Responsibility} \\
\midrule
Formatting & \modpath{formatting.py} & Pre-formatted plain text
             output using unicode box-drawing.  Used by both CLI
             and MCP for consistent output. \\
Sitemap    & \modpath{sitemap.py}    & Sitemap discovery and URL
             extraction for web ingestion. \\
Doctor     & \modpath{doctor.py} + \modpath{doctor\_*.py} & Health checks
             and install flow: model availability, database access, MCP
             configuration, daemon, sync, captures, ethos memory, recall
             telemetry (\modpath{doctor\_memory.py}, \modpath{doctor\_ethos.py},
             \modpath{doctor\_recall.py}). \\
Recall Telemetry & \modpath{query\_log.py} + \modpath{query\_log\_schema.py}
             & Local, engine-free SQLite store for recall telemetry
             (DES-056): one \texttt{query\_events} row + N
             \texttt{query\_hits} rows per search, written from
             \modpath{daemon/routes/search.py} and read back through
             \texttt{GET /insights}. Mirrors \modpath{sync\_registry.py}'s
             connection lifecycle. \\
Enable     & \modpath{enable.py}     & Project activation and
             deactivation: register directory, refresh the vendored
             ethos memory-guide blocks (\modpath{ethos\_ext\_block.py},
             \modpath{ethos\_ext\_scan.py}), write per-project config. \\
TLS        & \modpath{tls.py}        & TLS certificate generation:
             self-signed EC\,P-256 CA and server cert with
             atomic file writes and TOFU-safe CA reuse. \\
Proxy      & \modpath{proxy.py}      & \texttt{mcp-proxy} binary
             download, SHA256 verification, and
             platform-specific installation. \\
Logging    & \modpath{logging\_config.py} & Structured logging
             configuration: rotating file handler + stderr,
             \texttt{QUARRY\_LOG\_LEVEL} override. \\
\bottomrule
\end{tabular}
\end{center}

% ===================================================================
\section{Daemon Model}
% ===================================================================

\subsection{Why a Daemon}

The embedding model (snowflake-arctic-embed-m-v1.5, int8 quantized
ONNX) occupies $\sim$200\,MB of RAM and takes $\sim$250\,ms to load.
MCP's stdio transport model spawns a new process per session.  Without
a daemon, five Claude Code tabs means five copies of the model in
memory and five cold starts.

\texttt{quarryd} loads the model once and keeps it resident, respawned
by launchd \texttt{KeepAlive} / systemd \texttt{Restart=always}.

\subsection{Process Boundary and Loopback Authentication (v2.2)}
\label{sec:proc-boundary}

\texttt{quarryd} (\modpath{daemon/launcher.py}) is a dedicated entry
point --- the sole engine process, making the client/engine boundary a
hard \emph{process} split.  Only \modpath{daemon/} imports the engine;
the \texttt{quarry serve} subcommand is deleted.

At startup \texttt{quarryd} refuses a remote-reachable bind that carries
no operator key (an auto-generated token would be false security --- it is
unreadable by the remote clients that need it), and for a loopback bind
mints a 256-bit \texttt{serve.token} when none is supplied.  The daemon
writes that token mode-0600 beside \texttt{serve.port} (shared
\modpath{run\_dir.py}) via an atomic \texttt{os.open(0o600)} +
tmp-rename, and requires it on every request --- closing the
multi-user-host exposure where any local UID could reach the
unauthenticated daemon on \texttt{127.0.0.1}.  Loopback classification
(\modpath{net.py}, \texttt{LoopbackPolicy}) is \texttt{ipaddress}-based:
\texttt{localhost}/\texttt{::1}/\texttt{127.0.0.0/8} are loopback,
\texttt{0.0.0.0} and unresolved names fail closed.

Clients resolve the target through \texttt{ClientConfig}
(\modpath{client/config.py}): for a loopback URL the bearer is read
\emph{live} from \texttt{serve.token} (it rotates every daemon restart,
so a stored token is never trusted); for a remote URL the stored bearer
is kept.  A missing loopback token fails closed with an actionable error
rather than a bare 401.  \texttt{ClientConfig} feeds the one
\texttt{QuarryClient} that the CLI, the \texttt{quarry mcp} stdio server,
hooks, and the library all use; \texttt{RemoteClient} is deleted.

The daemon serves the \textbf{REST API only} under \texttt{/v1} (the
in-daemon MCP route was removed in PR-1); its request-admission trust model
combines the \texttt{JsonContentTypeGuard} --- an ASGI CSRF choke point
that rejects any \texttt{POST} without \texttt{Content-Type:
application/json} (a browser cannot send that cross-origin without a
preflight the CORS policy refuses) --- with the \texttt{serve.token} bearer
above.  The guard blocks a forged cross-origin write; the token blocks an
unauthorized local reader.

\subsection{Client/Engine Boundary (I1)}

The daemon-first split (DES-031) is a hard boundary, enforced
\emph{structurally} rather than by convention: a client process never
imports the engine.  Three layers, dependency arrow pointing inward:
\modpath{quarry/api} (the wire contract --- Pydantic models, zero engine,
zero heavy deps) $\leftarrow$ \modpath{quarry/client} (\texttt{QuarryClient}
transport + typed errors, imports only \texttt{api}) $\leftarrow$ the client
processes (\modpath{\_\_main\_\_}, \modpath{hooks}, \modpath{mcp\_server}).
The engine --- \modpath{db/}, \modpath{embeddings}, \modpath{ingestion/},
\modpath{retrieval/}, \modpath{sync}, \modpath{daemon/} --- is server-side
only; its entry point \modpath{daemon/launcher.py} (the \texttt{quarryd}
script) lives inside \modpath{daemon/}, so it is engine-side and never a
client.

Two mechanisms lock it (PR-6):

\begin{itemize}
  \item \textbf{An import-linter contract} (\texttt{.importlinter}, run by
    \texttt{make check-imports} in \texttt{make check} and CI) forbids any of
    the five client modules from importing any of the six engine packages
    along a \emph{transitive} chain; a violating import fails CI, not review.
    The one sanctioned exception is an explicit, enumerated
    \texttt{ignore\_imports} list of the host-admin diagnostic edges
    (\texttt{doctor}/\texttt{install}/\texttt{uninstall} probe the local engine
    environment to report why a daemon is unhealthy) --- function-body lazy
    imports off the hot path, not an invisibility trick (grimp resolves lazy
    imports).
  \item \textbf{A runtime engine-sabotage test} imports each client module in
    a subprocess with \texttt{lancedb}/\texttt{onnxruntime}/\texttt{pyarrow}
    poisoned and asserts no engine module reaches \emph{module} scope --- the
    companion guard that the diagnostic lazy imports stay deferred.
\end{itemize}

Daemon-dependent tests run against a hermetic in-process ASGI fixture
(\texttt{httpx.ASGITransport} over the FastAPI \texttt{build\_app}, no
socket); a single real-loopback-TLS smoke (slow tier, excluded from the fast
suite) covers the pinned-CA wire contract end-to-end; and \texttt{install}
gates on \texttt{/health} \texttt{ready} before exiting~0.

\subsection{MCP Transport (\texttt{quarry mcp})}

Quarry's plugin runs \texttt{quarry mcp} directly --- a FastMCP
\emph{stdio} server (\modpath{mcp\_server.py}) whose tool bodies
are thin \texttt{QuarryClient} calls over the daemon's \texttt{/v1} REST
API.  It reads MCP JSON-RPC from stdin, dispatches each tool to the
resident daemon (reusing its one warm \texttt{ctx.embedder}/\texttt{ctx.database}),
and relays the response to stdout.  It imports no engine, so
\texttt{import quarry.mcp\_server} and running \texttt{quarry mcp} load
zero LanceDB/ONNX.  Database selection is a per-request parameter carried
in the wire call; the daemon holds no per-connection \texttt{ContextVar}
because the in-daemon \texttt{/mcp} WebSocket route was removed (DES-031
v2.2, PR-1).  \texttt{mcp-proxy} is no longer on this path (it remains a
supported MCP stdio$\leftrightarrow$daemon transport for other MCP clients).

\subsection{CLI and the Daemon}

The CLI is a thin client: it builds a request, calls the daemon over the
\texttt{/v1} wire protocol via \texttt{QuarryClient}, and renders the
response --- it does not load the embedding model or open LanceDB
in-process.  The prior design had \texttt{quarry find "margins"} load the
model in-process and query LanceDB directly; that local-engine path is
removed.  The one sanctioned exception is the host-admin diagnostics
(\texttt{doctor}/\texttt{install}/\texttt{uninstall}), which probe the
local engine environment through function-body lazy imports so they can
report \emph{why} a daemon is unhealthy even when it is down --- an
explicit, enumerated \texttt{ignore\_imports} list, proven off the hot
path by the sabotage test.

% ===================================================================
\section{MCP Tool Surface}
% ===================================================================

\begin{center}
\small
\setlength{\tabcolsep}{4pt}
\begin{longtable}{lp{3.5cm}p{3.2cm}l}
\toprule
\textbf{Tool} & \textbf{Key Parameters} & \textbf{Purpose} &
  \textbf{Exec} \\
\midrule
\endfirsthead
\toprule
\textbf{Tool} & \textbf{Key Parameters} & \textbf{Purpose} &
  \textbf{Exec} \\
\midrule
\endhead
\tool{find}        & \field{query}, \field{limit},
  \field{collection}, \field{agent\_handle},
  \field{memory\_type}     & Hybrid search              & Sync \\
\tool{ingest}      & \field{source}, \field{overwrite},
  \field{collection}       & Ingest a URL (files go through
                             \tool{register\_directory}) & Bg \\
\tool{remember}    & \field{content},
  \field{document\_name}, \field{agent\_handle},
  \field{memory\_type}     & Ingest inline text         & Bg \\
\tool{learn}       & \field{lesson}, \field{topic},
  \field{name}             & Save a distilled lesson
                             (\texttt{memory\_type = lesson}) & Bg \\
\tool{list}        & \field{kind}
                           & List indexed data          & Sync \\
\tool{show}        & \field{document\_name},
  \field{page\_number}     & Document metadata or page  & Sync \\
\tool{delete}      & \field{name}, \field{kind}
                           & Delete document or
                             collection                 & Bg \\
\tool{register\_directory}   & \field{directory},
  \field{collection}       & Register for sync          & Bg \\
\tool{deregister\_directory} & \field{collection},
  \field{keep\_data}       & Remove registration        & Bg \\
\tool{sync\_all\_registrations} & ---  & Sync all registered dirs & Bg \\
\tool{status}      & ---  & Database stats             & Sync \\
\tool{use}         & \field{name}
                           & Switch active database     & Sync \\
\tool{missions\_sync} & \field{mission}, \field{dry\_run},
  \field{force}            & File frozen ethos mission rounds
                             into \texttt{memory-<worker>} & Sync$^{*}$ \\
\bottomrule
\end{longtable}
\end{center}

Bg = background (daemon-side; the tool returns a \field{task\_id}).
$^{*}$\tool{missions\_sync} runs client-side: it reads the repo's
\texttt{.punt-labs/ethos/missions/} tree and issues one
\texttt{/v1/remember} (202) per frozen round the daemon does not yet hold,
returning after the last post (DES-055).  The table above is the
registered surface: every tool on \texttt{McpTools} plus
\texttt{missions\_sync} on the sibling \texttt{MissionTools}
(\modpath{mcp\_missions.py}), registered from one call; the exact set is
pinned by \texttt{tests/test\_mcp\_server.py}.

Each tool body is a thin \texttt{QuarryClient} call over \texttt{/v1}
(the MCP server holds no engine).  A background tool posts to the daemon,
which returns \texttt{202~Accepted} with a \field{task\_id}; the actual
work runs on the daemon's serialized per-\texttt{RouteKey} queue (DES-042)
behind the embed \texttt{Semaphore}, and the caller polls
\texttt{/v1/tasks/\{task\_id\}} or ignores the outcome.  This keeps a
30-second ingest from blocking the LLM response stream while preventing
the concurrent-write races the single-writer-per-table invariant closes.
Read tools return their result inline.

% ===================================================================
\section{HTTP Server}
% ===================================================================

\subsection{ASGI Application}

FastAPI app built by \texttt{build\_app()} (\modpath{daemon/app.py}),
served by uvicorn.  FastAPI \emph{is} Starlette underneath, so request
validation \emph{and} the OpenAPI document derive from the same
\modpath{quarry/api} Pydantic models --- no hand-rolled emitter or route
registry.  REST handlers are sync functions run in the threadpool; the
routes live in \modpath{daemon/routes/}.  There is no WebSocket endpoint
--- the daemon serves the \textbf{REST API only} (the in-daemon
\texttt{/mcp} route was removed in DES-031 v2.2, PR-1).

Default port: 8420 (\texttt{DEFAULT\_PORT} in \modpath{config.py}),
overridable with \texttt{-{}-port}.  All engine operations sit under
\texttt{/v1}; \texttt{/health} and \texttt{/ca.crt} are unversioned.
\texttt{/health} advertises \field{api\_version}; a major mismatch
raises \texttt{QuarryVersionError} in the client.

\subsection{Endpoints}

\begin{center}
\begin{tabular}{llp{5.3cm}ll}
\toprule
\textbf{Path} & \textbf{Method} & \textbf{Purpose} &
  \textbf{Auth} & \textbf{Model} \\
\midrule
\texttt{/health}              & GET    & Server state + \field{api\_version} & No  & sync  \\
\texttt{/ca.crt}              & GET    & CA certificate PEM (TOFU)        & No  & sync  \\
\texttt{/v1/search}           & GET    & Hybrid search                    & Yes & sync  \\
\texttt{/v1/show}             & GET    & Document metadata or page text   & Yes & sync  \\
\texttt{/v1/documents}        & GET    & List documents                   & Yes & sync  \\
\texttt{/v1/documents}        & DELETE & Delete document (returns 202)    & Yes & async \\
\texttt{/v1/collections}      & GET    & List collections                 & Yes & sync  \\
\texttt{/v1/collections}      & DELETE & Delete collection (returns 202)  & Yes & async \\
\texttt{/v1/registrations}    & GET    & List directory registrations     & Yes & sync  \\
\texttt{/v1/registrations}    & POST   & Register directory (returns 202) & Yes & async \\
\texttt{/v1/registrations}    & DELETE & Deregister directory (sync + 202 purge) & Yes & async \\
\texttt{/v1/databases}        & GET    & List named databases             & Yes & sync  \\
\texttt{/v1/status}           & GET    & Database statistics              & Yes & sync  \\
\texttt{/v1/insights}         & GET    & Recall telemetry snapshot (DES-056) & Yes & sync  \\
\texttt{/v1/coverage}         & GET    & Per-collection coverage counts   & Yes & sync  \\
\texttt{/v1/use}              & POST   & Switch active database           & Yes & sync  \\
\texttt{/v1/remember}         & POST   & Ingest inline content (returns 202) & Yes & async \\
\texttt{/v1/learn}            & POST   & Save a distilled lesson (returns 202) & Yes & async \\
\texttt{/v1/ingest}           & POST   & Ingest a URL (returns 202)       & Yes & async \\
\texttt{/v1/capture}          & POST   & Capture web/transcript (returns 202) & Yes & async \\
\texttt{/v1/captures/lookup}  & POST   & Prior capture of a URL (web-fetch loop closer) & Yes & sync  \\
\texttt{/v1/captures/push}    & POST   & Push private capture shadow repo & Yes & async \\
\texttt{/v1/sync}             & POST   & Sync registrations (returns 202) & Yes & async \\
\texttt{/v1/optimize}         & POST   & Compact + rebuild index (returns 202) & Yes & async \\
\texttt{/v1/backfill-sessions}& POST   & Backfill session transcripts (returns 202) & Yes & async \\
\texttt{/v1/tasks/\{task\_id\}} & GET  & Poll task status (unified)       & Yes & sync  \\
\bottomrule
\end{tabular}
\end{center}

\texttt{/v1/tasks/\{task\_id\}} is the single, unified poll endpoint ---
the old \texttt{/sync/\{id\}} and \texttt{/ingest/\{id\}} alias routes
were removed in the DES-031 v2.2 contract rewrite.  See
Appendix~\ref{app:concurrency} for the full concurrency matrix across
HTTP, CLI, and MCP surfaces.

\subsection{Port File}

After binding, the daemon writes the actual port to
\texttt{\~{}/.punt-labs/quarry/data/<db>/serve.port}, and (when it has an
effective key) a mode-0600 \texttt{serve.token} beside it (see
\S\ref{sec:proc-boundary}).  This supports \texttt{-{}-port 0} (ephemeral
assignment).  Removed on shutdown via the lifespan context.

% ===================================================================
\section{Configuration}
% ===================================================================

\subsection{Settings}

All settings are environment variables read via pydantic-settings.

\begin{center}
\small
\setlength{\tabcolsep}{4pt}
\begin{longtable}{llp{6.5cm}}
\toprule
\textbf{Variable} & \textbf{Default} & \textbf{Description} \\
\midrule
\endfirsthead
\toprule
\textbf{Variable} & \textbf{Default} & \textbf{Description} \\
\midrule
\endhead
\multicolumn{3}{l}{\textbf{Paths}} \\
\midrule
\texttt{QUARRY\_ROOT}    & \texttt{\~{}/.punt-labs/quarry/data} &
  Base directory for all databases \\
\texttt{LANCEDB\_PATH}   & \texttt{<root>/default/lancedb} &
  Vector database location \\
\texttt{REGISTRY\_PATH}  & \texttt{<root>/default/registry.db} &
  Directory sync SQLite database \\
\texttt{TELEMETRY\_PATH} & \texttt{<root>/default/telemetry.db} &
  Recall-telemetry SQLite database (DES-056) \\
\texttt{QUARRY\_LOG\_DIR} & \texttt{\~{}/.punt-labs/quarry/logs/} &
  Log directory (\texttt{quarry.log} for the CLI/MCP client tier,
  \texttt{quarryd.log} for the daemon --- DES-048; rotating, 5\,MB
  $\times$ 5; resolvable per call via \texttt{QUARRY\_LOG\_DIR}) \\
\midrule
\multicolumn{3}{l}{\textbf{Chunking}} \\
\midrule
\texttt{CHUNK\_MAX\_CHARS}     & \texttt{1800} &
  Target max characters per chunk ($\sim$450 tokens) \\
\texttt{CHUNK\_OVERLAP\_CHARS} & \texttt{200} &
  Overlap between consecutive chunks \\
\midrule
\multicolumn{3}{l}{\textbf{Embedding}} \\
\midrule
\texttt{EMBEDDING\_MODEL}     & \texttt{Snowflake/\allowbreak snowflake-arctic-embed-m-v1.5} &
  Model identifier (display only --- the ONNX model is fixed) \\
\texttt{EMBEDDING\_DIMENSION} & \texttt{768} &
  Vector dimension (display only --- fixed at 768) \\
\midrule
\multicolumn{3}{l}{\textbf{Progressive commit (DES-034)}} \\
\midrule
\texttt{SYNC\_FLUSH\_MB}       & \texttt{32} &
  Flush buffered vectors to LanceDB once they reach this many MB
  (bounds peak resident vectors; may fire mid-file) \\
\texttt{EMBED\_WINDOW\_CHUNKS} & \texttt{512} &
  Chunks embedded per window, so a large document never materializes
  all its vectors (a kpz throughput seam) \\
\midrule
\multicolumn{3}{l}{\textbf{Recall telemetry (DES-056)}} \\
\midrule
\texttt{TELEMETRY\_ENABLED}         & \texttt{true} &
  Record one \texttt{query\_events} row + N \texttt{query\_hits} rows per
  search; \texttt{false} makes the write path a true no-op \\
\texttt{TELEMETRY\_RETENTION\_DAYS} & \texttt{90} &
  Prune query-log rows older than this many days \\
\midrule
\multicolumn{3}{l}{\textbf{Watch/index loop (DES-045)}} \\
\midrule
\texttt{WATCH\_ENABLED}        & \texttt{true} &
  Run the always-on watch/index loop (installing \texttt{quarryd} is the
  opt-in; reverses the prior timer posture) \\
\texttt{WATCH\_DEBOUNCE\_S}    & \texttt{1.0} &
  Coalesce an edit burst before submitting \\
\texttt{WATCH\_MAX\_DELAY\_S}  & \texttt{5.0} &
  Cap on a continuously-rearmed path (anti-starvation) \\
\texttt{WATCH\_BULK\_THRESHOLD} & \texttt{50} &
  A delta above this many paths collapses to one bulk
  \texttt{CollectionSyncJob} \\
\texttt{WATCH\_USE\_POLLING}   & \texttt{false} &
  Force watchdog's stat-walk \texttt{PollingObserver} \\
\texttt{WATCH\_SAFETY\_SCAN\_S} & \texttt{300.0} &
  Periodic roster reconcile + orphan sweep; picks up new databases and
  re-submits shed scans (0 = off) \\
\texttt{WATCH\_OPTIMIZE\_MIN\_INTERVAL\_S} & \texttt{30.0} &
  \texttt{FinalizeThrottle}: at most one finalize (optimize + FTS
  rebuild) per interval \emph{per database} (0 = every batch) \\
\midrule
\multicolumn{3}{l}{\textbf{Thread limits (process env; DES-032/045c)}} \\
\midrule
\texttt{LANCE\_CPU\_THREADS} / \texttt{LANCE\_IO\_THREADS} & \texttt{2} &
  Clamp LanceDB's tokio compute/IO pools (fail-closed, floor == cap == 2;
  set before the first \texttt{lancedb.connect}) \\
\texttt{OMP\_NUM\_THREADS}     & \texttt{min(2, ncpu)} &
  Cap OpenMP to match the ONNX intra-op limit (floor 1) \\
\bottomrule
\end{longtable}
\end{center}

The LanceDB and OMP caps are applied by \texttt{ThreadConfig}
(\modpath{thread\_config.py}) in \texttt{DaemonServer.run()} via a
fail-closed clamp: an inherited or operator value is honored only within
\texttt{[floor, cap]}; a higher or non-numeric value is forced \emph{down}
to the cap (a \texttt{setdefault} would yield upward to a stale export and
void the ceiling) and a below-floor value is raised to the floor.  ONNX
\texttt{intra\_op\_num\_threads} is \texttt{min(2, ncpu)} on CPU and~1 on
GPU; \texttt{inter\_op\_num\_threads} is~1 (DES-032).

\subsection{Named Databases}

\texttt{resolve\_db\_paths(settings, name)} maps a database name to a
path under \texttt{QUARRY\_ROOT}.  Each database is fully isolated ---
its own LanceDB directory, sync registry, and vector index.  Names containing path separators (\texttt{/} or
\texttt{\textbackslash}) or the literal traversal segments \texttt{.}
or \texttt{..} are rejected; names like \texttt{foo.bar} are allowed.
If \texttt{LANCEDB\_PATH} was explicitly overridden, the override is
preserved.

The persistent default database is stored in
\texttt{\~{}/.punt-labs/quarry/config.toml} as \texttt{[default] database = "name"}.

\subsection{Chunking Tuning}

\textbf{Increase \texttt{CHUNK\_MAX\_CHARS}} (3000--4000) for documents
with long self-contained sections: legal contracts, academic papers,
large source code functions.

\textbf{Decrease \texttt{CHUNK\_MAX\_CHARS}} (800--1200) for dense
documents where each paragraph covers a distinct topic.

\textbf{Increase \texttt{CHUNK\_OVERLAP\_CHARS}} (400) when relevant
information spans paragraph boundaries.

The embedding model has a 512-token context window.  Chunks exceeding
this are truncated during embedding, reducing quality for the tail.

% ===================================================================
\section{Search and Retrieval}
% ===================================================================

\subsection{Embedding Model}

snowflake-arctic-embed-m-v1.5 via ONNX Runtime:

\begin{itemize}
  \item 768-dimensional vectors, 512-token context window
  \item Asymmetric retrieval: queries are prefixed with a search
    instruction, documents are not
  \item Output vectors are L2-normalized to unit length, so search runs
    under the \textbf{cosine} metric and similarity
    ($1 - \text{distance}$) is bounded true cosine in $[-1, 1]$
    (DES-038)
  \item int8 quantized ONNX for CPU ($\sim$120\,MB), FP16 for CUDA
    ($\sim$220\,MB)
  \item Provider auto-detection: CUDA+FP16 when available, CPU+int8
    fallback.  \texttt{QUARRY\_PROVIDER} env var overrides
    (\texttt{cpu}, \texttt{cuda}, or unset for auto).
    See DES-016 and \modpath{ingestion/provider.py}.
  \item Auto-downloads on first use via \texttt{huggingface\_hub}
\end{itemize}

The model is fixed --- there is no configuration to swap it.  Changing
the model invalidates all existing embeddings, requiring full
re-ingestion.

\subsection{OCR}

RapidOCR handles all image-based text extraction.  For semantic
search, OCR accuracy matters less than expected --- the embedding
model is robust to minor OCR errors.  A misspelled word rarely
changes semantic meaning enough to affect search ranking.

\subsection{Search Quality Tips}

\begin{enumerate}
  \item \textbf{Use natural language queries.}  The model is trained on
    question-passage pairs.  ``What were Q3 revenue figures?'' works
    better than ``Q3 revenue.''
  \item \textbf{Use collection filtering.}
    \texttt{quarry find "auth logic" -c backend}.
  \item \textbf{Re-ingest after tuning.}  Chunk parameters only affect
    new ingestions.  Run \texttt{quarry sync} after changes.
  \item \textbf{Check chunk boundaries.}
    \texttt{quarry find -n 1 "query"} and examine the result.
  \item \textbf{Use \texttt{-{}-overwrite} when re-ingesting.}
    Without it, duplicate chunks accumulate.
\end{enumerate}

\subsection{Hybrid Search}

Hybrid search is the \textbf{default for every query} --- there is no
vector-only path.  All three surfaces (CLI, HTTP, MCP) construct a
\texttt{SearchService} (\modpath{retrieval/service.py}) and call its one
\texttt{search} method, so the surfaces cannot drift into different
retrieval behavior.  This closes bug class 3 (remote/local divergence),
where the HTTP \texttt{/search} endpoint once ran vector-only search
while the CLI ran hybrid.  \texttt{SearchService} wraps a
\texttt{HybridRetriever} parameterized by a \texttt{RetrievalConfig}
whose defaults reproduce the shipped pipeline bit-for-bit.

\textbf{Channels:}
\begin{enumerate}
  \item \textbf{Vector similarity} --- ANN search via LanceDB under the
    cosine metric
  \item \textbf{BM25 full-text search} --- LanceDB's \emph{native}
    (Rust) inverted index on the \texttt{text} column, the FTS backend
    LanceDB builds by default, created during schema migration.  Quarry's
    \texttt{LanceTable} protocol and the \texttt{RecyclingTable} wrapper
    forward only \texttt{replace} to \texttt{create\_fts\_index}; there is
    no separate Tantivy \texttt{IndexWriter} process or thread pool, so the
    index rides the daemon's capped LanceDB runtime (DES-045c).  (Earlier
    ADRs --- DES-017, DES-023 --- name the backend ``Tantivy''; that name is
    legacy, the shipped engine is LanceDB-native.)
\end{enumerate}

Results are fused using Reciprocal Rank Fusion (RRF):
\[
  \text{score}(m) = \sum_{\text{channel}} \frac{1}{60 + \text{rank}_{\text{channel}}(m)}
\]

\textbf{FTS scoring.}  FTS rows arrive from the full-text channel without
a vector distance.  Rather than displaying a synthetic $1.00$ (which made
irrelevant keyword hits appear as perfect matches), each FTS-only row is
annotated with its \emph{true} cosine distance against the query vector;
a row whose stored vector is missing or zero-length is assigned the
worst-case distance (similarity $-1$) instead.

\textbf{Temporal decay} is not a second retrieval channel --- it is an
additional weighting term applied only to agent-scoped memories (non-empty
\texttt{agent\_handle}).  An \texttt{agent\_handle} filter therefore
narrows the result set and enables decay; it does not switch retrieval
modes.
\[
  w = e^{-\lambda \cdot h}
\]
where $\lambda$ is the decay rate and $h$ is hours since ingestion.
Unscoped documents ($\texttt{agent\_handle} = \text{``''}$) are exempt
from decay (production leaves $\lambda = 0$).

\subsection{Retrieval Seam and Evaluation}

The \modpath{retrieval/} package is a deliberate seam (DES-037): the same
\texttt{HybridRetriever} that serves production also serves the retrieval
evaluation harness (\texttt{make eval}), so the committed baseline
measures the retriever that actually ships.  \texttt{RetrievalConfig}
exposes the levers a bake-off varies --- RRF $k$, over-fetch multiplier,
distance metric, exact-vs-ANN vector search, reranker, and embedding
strategy --- without forking the code path.  The harness is internal
dev/CI tooling; it is not an end-user command.  See
\texttt{docs/eval-harness-design.md}.

Three columns extend the chunks table for agent-scoped memory:

\begin{center}
\begin{tabular}{llll}
\textbf{Column} & \textbf{Type} & \textbf{Default} & \textbf{Purpose} \\
\hline
\texttt{agent\_handle} & utf8 & \texttt{""} & Agent identity scope \\
\texttt{memory\_type} & utf8 & \texttt{""} & fact, observation, opinion, procedure, lesson \\
\texttt{summary} & utf8 & \texttt{""} & One-line summary \\
\end{tabular}
\end{center}

Migration is idempotent --- \texttt{\_migrate\_schema()} adds missing columns
with empty-string defaults on every table open. Existing databases gain
the columns transparently.

\texttt{memory\_type} is a closed vocabulary (\modpath{memory\_types.py},
\texttt{MemoryType}): the four agent-owned types decay; \texttt{lesson} is
written only by \texttt{learn} and is boosted instead of decayed.  Every
write route --- \texttt{/v1/remember}, \texttt{/v1/ingest},
\texttt{/v1/capture} --- validates through one \texttt{RouteGroup} method,
so an unknown value is a 400 with an identical body on all three, never a
silently stored row that neither decays nor matches a typed filter
(DES-055).

Identity tagging: the parent-session hooks (PreCompact, SessionEnd) read
the agent handle from the ethos repo pin --- \texttt{.punt-labs/ethos.yaml}
first, then the legacy \texttt{.punt-labs/ethos/config.yaml}, at each
ancestor up to the repository boundary (\modpath{ethos\_tree.py},
\modpath{ethos\_handle.py}) --- and tag captured transcripts with it.
SubagentStop never uses the pin for attribution (the pin names the
leader): it uses \texttt{agent\_type} iff that handle is a registered
identity in the vendored or global ethos tree, else the capture is filed
unattributed.  A \texttt{remember} carries the caller's own handle; the
daemon never infers one from \texttt{cwd} (DES-055).

\subsection{Ethos Extension Integration}

Quarry stores its configuration in ethos's generic extension mechanism
(DES-008, DES-019):

\begin{verbatim}
~/.punt-labs/ethos/identities/<handle>.ext/quarry.yaml   (global)
.punt-labs/ethos/identities/<handle>.ext/quarry.yaml    (vendored, per repo)
\end{verbatim}

The extension file contains two categories of data:

\begin{enumerate}
  \item \textbf{Sidecar config} --- \texttt{memory\_collection} and
    \texttt{expertise\_collections} are read directly by quarry via the
    filesystem (no ethos API call).
  \item \textbf{Session context} --- \texttt{session\_context} is a YAML
    literal block scalar emitted verbatim by ethos at session start and
    before context compaction.  It contains the agent's memory instructions
    (collection name, slash commands, memory types).
\end{enumerate}

The \texttt{session\_context} carries a versioned \texttt{\#\# Memory
(quarry guide v2)} block (\modpath{ethos\_ext\_block.py}): the five
moments to \texttt{remember}, what never to store, and the rule that an
agent always passes its own handle.  \texttt{quarry install} refreshes the
global identities; \texttt{quarry enable} refreshes the vendored
\texttt{.punt-labs/ethos/identities/<handle>.ext/quarry.yaml} files (a
committed diff, reported as ``commit via PR''); \texttt{doctor} stays
read-only.  A v1 block is replaced in place, a customised block is left
alone, and per-identity outcomes (\modpath{ethos\_ext\_scan.py}) ensure one
malformed file does not abort the rest.

Ownership boundary: ethos emits quarry's instructions without understanding
them.  Quarry writes its own instructions without importing ethos.  The
dependency is one-way --- quarry depends on ethos for identity; ethos has
zero knowledge of quarry's internals.

\subsection{Agent-Memory Write Loop (DES-055)}

Recall-on-launch had worked since the agent-memory schema landed: ethos
injects each identity's \texttt{session\_context} into a spawned agent.
The \emph{write} side stayed dormant --- no \texttt{memory-*} collection
existed --- until DES-055 added the three producers that fill it.  None of
them runs a model; each is a thin client of the daemon's
scrub-before-store \texttt{/v1/remember} route.

\begin{enumerate}[label=\textbf{Loop \arabic*}.]
  \item \textbf{Persist habit.}  The memory guide above tells every agent
    when to \texttt{remember} and to pass its own handle.  The daemon
    routes a handle-bearing \texttt{remember} with an empty
    \texttt{collection} to \texttt{memory-<handle>}.
  \item \textbf{Evaluator feedback.}  \texttt{quarry missions sync} (CLI),
    \texttt{missions\_sync} (MCP), and \texttt{/quarry missions sync}
    (plugin) read the repo's \texttt{.punt-labs/ethos/missions/} trio
    (\texttt{contract}, \texttt{results}, \texttt{reflections}) and file
    each \emph{frozen} round as an \texttt{observation} in
    \texttt{memory-<worker>}, named
    \texttt{mission-<repo>-<id>-r<n>}.  Idempotent: a round the daemon
    already holds is skipped, a name held by another checkout's round is
    an error and never overwritten (\texttt{-{}-force} re-files a matching
    key only), \texttt{-{}-dry-run} posts nothing.  Ethos is never called,
    so the one-way dependency holds.
  \item \textbf{Capture distillation.}  \texttt{SubagentStop} archives the
    subagent's own transcript to \texttt{<repo>-captures} and, when
    \texttt{agent\_type} names a registered identity, files the subagent's
    final assistant turn --- already a summary by construction --- as an
    \texttt{observation} in \texttt{memory-<handle>}
    (\texttt{subagent-<id8>-report}).  A bare \texttt{Agent()}
    (\texttt{general-purpose}) is filed unattributed.
\end{enumerate}

\begin{center}
\begin{tabular}{llp{6.6cm}}
\toprule
\textbf{Module} & \textbf{File} & \textbf{Responsibility} \\
\midrule
Vocabulary  & \modpath{memory\_types.py}      & \texttt{MemoryType} ---
              the closed \texttt{memory\_type} vocabulary and its
              write-rejection rule (one source for routes, decay, doctor,
              and CLI help). \\
Ethos tree  & \modpath{ethos\_tree.py},
              \modpath{ethos\_handle.py}      & Read-only locator for the
              ethos sidecar (pins, vendored/global identities, missions);
              pin-based and \texttt{agent\_type}-based handle resolution. \\
Mission store & \modpath{mission\_store.py},
              \modpath{mission\_records.py},
              \modpath{mission\_round\_parts.py} & Read every mission under
              the repo's missions tree into \texttt{MissionContract} +
              \texttt{MissionRound} records; parse errors are collected,
              never fatal. \\
Mission memory & \modpath{mission\_memory.py},
              \modpath{mission\_sync\_types.py} & \texttt{MissionMemorySync}
              --- compose one \texttt{RememberRequest} per frozen round,
              check existence via \texttt{show}, post, tally
              filed/skipped/errors. \\
Surfaces    & \modpath{cli\_missions.py},
              \modpath{mcp\_missions.py}      & The \texttt{missions sync}
              Typer sub-app and the \texttt{missions\_sync} MCP tool; both
              drive the same \texttt{MissionMemorySync}. \\
Subagent    & \modpath{subagent\_capture.py},
              \modpath{subagent\_report.py}   & Two rows from one event:
              raw transcript capture plus the distilled report. \\
\bottomrule
\end{tabular}
\end{center}

\subsection{Recall Telemetry (DES-056)}

DES-055 gave agents a place to write memory; DES-056 answers whether recall
is actually exercised.  \modpath{daemon/routes/search.py} times every
\texttt{SearchService.search()} call and, when
\texttt{Settings.telemetry\_enabled}, scrubs the query text through the same
\texttt{Scrubber} pass captures use (DES-041) and records one
\texttt{query\_events} row plus one \texttt{query\_hits} row per ranked
result to a local, engine-free SQLite store (\modpath{query\_log.py} +
\modpath{query\_log\_schema.py}), mirroring \modpath{sync\_registry.py}'s
WAL/busy-timeout connection lifecycle.  The write is boundary I/O: wrapped in
\texttt{try/except}, a failure is logged and swallowed so a telemetry outage
never turns a search into a 500.

\texttt{GET /v1/insights} mirrors the \texttt{status} chain across all three
surfaces (\texttt{quarry insights}, the \texttt{insights} MCP tool, the HTTP
route) and reports total queries, empty-result rate, p50/p95 latency, top
empty (scrubbed) queries, per-collection and per-agent recall counts, a
memory-vs-knowledge split, and a hit recency (decay-band) breakdown.  The
decay-band join --- a hit's document identity against the live catalog's
\texttt{ingestion\_timestamp} --- runs in the route, not the store: the
query log has no LanceDB import, so only the code already holding
\texttt{ctx.database} can perform it.  \texttt{SearchRequest} carries an
optional \texttt{surface} field (\texttt{cli}/\texttt{mcp}/\texttt{http}/
\texttt{plugin}) for provenance; the CLI \texttt{find} path stamps
\texttt{surface="cli"} the same way the MCP \texttt{find} tool stamps
\texttt{"mcp"}.

\subsection{Skills and Cross-Harness Distribution (quarry-x7ja)}

Quarry ships two Agent Skills at \texttt{plugin/skills/}: \texttt{quarry-recall}
(retrieve --- understand code, recall a project decision, recall memory, all
backed by \texttt{quarry find} with different scope flags) and
\texttt{quarry-capture} (persist --- \texttt{remember}/\texttt{learn}/
\texttt{ingest}, the five-moments timing, the \texttt{agent\_handle}
discipline).  Both drive the \texttt{quarry} CLI directly rather than naming
MCP tools, so the same skill content works unmodified on any harness; deep
detail (the failure-mode recovery table, scoping and decay, the memory-type
vocabulary) lives in each skill's \texttt{references/} for progressive
disclosure, while the SessionStart trigger trailer, the MCP server's
\texttt{instructions} block, and the ethos-injected memory guide are kept to
short pointer sentences naming the two skills, rather than each independently
restating the same how-to.

Claude Code reads \texttt{plugin/skills/} directly through the marketplace
plugin tree --- no deposit step.  \modpath{skills\_install.py}
(\texttt{SkillsInstaller}) composes \modpath{skills\_source.py}
(\texttt{SkillSource}, which owns every source-tree and single-deposit
mechanic: locate, hash, manifest, the atomic copy-and-swap) and adds only
harness orchestration on top --- which harness, which root. Together they
copy the same two skill directories, verbatim, into the auto-load location
each other harness was confirmed (not guessed) to read skills from:

\begin{itemize}
  \item \textbf{pi} --- \texttt{\textasciitilde{}/.pi/agent/skills/<name>/}, matching
    \texttt{getAgentDir()}'s resolution to \texttt{\textasciitilde{}/.pi/agent}.
  \item \textbf{opencode} ---
    \texttt{\textasciitilde{}/.config/opencode/skills/<name>/},
    one of the paths opencode's own skill loader globs
    (\texttt{\{skill,skills\}/**/SKILL.md}) as its ``global skills'' root.
  \item \textbf{codex} --- \texttt{\textasciitilde{}/.codex/skills/<name>/} (or
    \texttt{\$CODEX\_HOME}), the location codex's own preinstalled
    \texttt{skill-installer} skill deposits third-party skills into,
    explicitly alongside the sibling \texttt{.system/} directory that holds
    codex's own preinstalled skills --- never inside it.
\end{itemize}

Each deposit is idempotent and version-stamped: a \texttt{.quarry-skill.json}
manifest beside the deposited \texttt{SKILL.md} records a short content hash,
so a re-run with unchanged source content is a no-op and a changed source
overwrites the stale copy wholesale.  The copy itself is swapped in
atomically --- written into a temp sibling directory first, the previous
deposit (if any) renamed to a backup sibling, then the temp directory renamed
into place, with the backup restored on any rename failure --- so an
interrupted write never leaves a half-deposited skill on disk (the file-I/O
safety class this repo's review history flags first).

\texttt{quarry skills install [--agent <id>|--all]} and \texttt{quarry skills
status} (\modpath{cli\_skills.py}) expose this: \texttt{--all} (the default)
targets every harness \texttt{SkillsInstaller.detected\_harnesses()} finds on
the machine (a harness's own top-level config directory, e.g. \texttt{\textasciitilde{}/.codex},
existing); \texttt{--agent} targets one harness explicitly, detected or not.

\texttt{quarry disable} retracts the deposits again through a private
\texttt{\_SkillRetraction} collaborator, gated by two guards, both required
(defense-in-depth for an \texttt{rmtree} under the operator's real
\texttt{\$HOME}): \textbf{identity} --- every marketplace-layout Claude Code
plugin (lux, prfaq, dungeon, punt-kit, quarry itself) ships an
identically-shaped \texttt{plugin/skills/} tree, so that shape alone proves
nothing; only a checkout whose own \texttt{pyproject.toml} names the
\texttt{punt-quarry} package is quarry's own, checked on the disabled repo's
root directly, never via an ancestor walk (an ancestor walk would
true-positive on every path nested inside the quarry checkout, including
every test's own temp directory, and reach the operator's real home
directory as a side effect of an unrelated \texttt{disable} call); and
\textbf{manifest ownership} --- \texttt{SkillsInstaller.remove}/
\texttt{remove\_all} refuse (\texttt{action="foreign"}) any same-named
directory that lacks a \texttt{.quarry-skill.json} manifest, so a directory
this installer never deposited is never an \texttt{rmtree} candidate no
matter how it got there. Retraction is best-effort: a filesystem failure
during either guard's checks or the removal itself is logged and swallowed,
never aborting \texttt{disable} or blocking the marker/import teardown that
already committed.

% ===================================================================
\section{Logging and Exceptions}
% ===================================================================

\subsection{Logging Standard}

Every module that performs I/O declares
\texttt{logger = logging.getLogger(\_\_name\_\_)}.  Pure data modules
(\texttt{models.py}, \texttt{types.py}, \texttt{chunker.py}) do not.

Logging is configured once at each application entry point --- the CLI,
the MCP server, and the \texttt{quarryd} daemon (DES-048; the daemon writes
its own \texttt{quarryd.log} with crash hooks for all three uncaught-exception
routes).  Library modules never call \texttt{logging.basicConfig}.

\begin{center}
\begin{tabular}{lp{8cm}}
\toprule
\textbf{Level} & \textbf{Use} \\
\midrule
\texttt{DEBUG}   & Internal state: variable values, branch decisions,
                   intermediate results. \\
\texttt{INFO}    & Operational milestones: starting a job, completing
                   a step, resource counts. \\
\texttt{WARNING} & Unexpected but recoverable: empty input, deprecated
                   parameter usage. \\
\texttt{ERROR}   & Failures that prevent completion but do not crash.
                   Always with exception context. \\
\bottomrule
\end{tabular}
\end{center}

\texttt{CRITICAL} is not used.  If the process must exit, it raises an
exception.

Log messages use \texttt{\%s}-style formatting (lazy evaluation), not
f-strings.  Each message includes enough context to trace an operation
without reading code.

\subsection{Exception Handling}

The codebase uses built-in exception types.  No custom exceptions
unless a caller needs to distinguish failure modes programmatically.

\begin{center}
\begin{tabular}{lp{8cm}}
\toprule
\textbf{Exception} & \textbf{Raised when} \\
\midrule
\texttt{ValueError}       & Invalid input: unsupported format, bad
                            parameter. \\
\texttt{FileNotFoundError} & A file path argument does not exist. \\
\texttt{RuntimeError}      & An operation reports failure (OCR engine
                            error, ingestion failure). \\
\texttt{TimeoutError}      & An external service exceeds its time
                            budget. \\
\bottomrule
\end{tabular}
\end{center}

\textbf{Rules:}

\begin{enumerate}
  \item Never swallow exceptions.  Every \texttt{except} block must
    re-raise, or log with \texttt{logger.exception()} and raise a
    different exception, or (at system boundaries only) return a
    documented sentinel.
  \item Never use bare \texttt{except:}.
  \item Never catch \texttt{Exception} broadly in library code.  Broad
    catches are permitted only at the outermost boundary (MCP tool
    handler, CLI command).
  \item Exceptions carry context: include the values that caused the
    failure.
  \item Use \texttt{try/finally} for cleanup, not
    \texttt{try/except}.
  \item Do not use exceptions for flow control.
\end{enumerate}

% ===================================================================
\section{Security}
% ===================================================================

\subsection{Bearer Token Authentication}

\texttt{quarryd -{}-api-key} (or \texttt{QUARRY\_API\_KEY} env var)
gates all HTTP endpoints except \texttt{/health} behind
\texttt{Authorization: Bearer <key>}.

\begin{itemize}
  \item \texttt{/health} is exempt --- load balancers need
    unauthenticated health checks.
  \item \texttt{OPTIONS} is exempt --- CORS preflight requests carry
    no auth.
  \item \texttt{hmac.compare\_digest} for token comparison ---
    prevents timing attacks.
  \item Case-insensitive scheme per RFC~7235.
  \item Empty key = no auth --- prevents accidental trivial bypass.
  \item No key = auth disabled --- local use requires no key.
\end{itemize}

The server refuses to bind to a non-loopback address without
\texttt{-{}-api-key}.

\subsection{Request Admission (CSRF)}

The daemon serves REST only, so there is no WebSocket to hijack; the
cross-origin threat is a forged browser \texttt{POST}.  Two checks
compose the admission trust model:

\begin{itemize}
  \item \texttt{JsonContentTypeGuard} --- an ASGI choke point that
    rejects any \texttt{POST} without \texttt{Content-Type:
    application/json}.  A browser cannot send that header cross-origin
    without a preflight the CORS policy refuses, so a forged
    cross-origin write is blocked at the door.
  \item The \texttt{serve.token} / API-key bearer (checked on every
    \texttt{/v1} request) blocks an unauthorized \emph{local} reader.
\end{itemize}

\subsection{CORS}

\texttt{CORSMiddleware} configured with GET and OPTIONS methods,
Authorization and Content-Type headers.  Default allowed origin:
\texttt{http://localhost}.  Configurable via \texttt{-{}-cors-origin}
(repeatable).

The server reflects the request's \texttt{Origin} header only when it
matches the allow list, and adds \texttt{Vary: Origin} to signal
cacheability varies by origin.  When no \texttt{Origin} is present,
no CORS headers are emitted.

\subsection{TLS and TOFU Certificate Pinning}

The installed service always runs with TLS enabled.  \texttt{quarryd}
without \texttt{-{}-tls} is supported for local development but
must not be used for production.  Certificate management uses TOFU
(Trust On First Use) with self-signed CA pinning (DES-020).

\textbf{Certificate generation:} \texttt{quarry install} generates
an EC~P-256 CA and server certificate.  Files are written atomically
to \texttt{\~{}/.punt-labs/quarry/tls/} with restricted permissions
(0600 for keys, 0644 for CA cert).  The server cert SAN includes the
configured hostname (via \texttt{QUARRY\_TLS\_HOSTNAME}), localhost,
and loopback addresses.

\textbf{TOFU login:}
\begin{enumerate}
  \item Client fetches \texttt{/ca.crt} over HTTPS (verification
    disabled --- bootstrap).
  \item Displays SHA256 fingerprint for out-of-band confirmation.
  \item Pins CA cert at
    \texttt{\~{}/.punt-labs/mcp-proxy/quarry-ca.crt}.
  \item All subsequent connections verify against the pinned CA only
    --- system roots excluded.
\end{enumerate}

\textbf{Pinned-CA consumers:} \texttt{QuarryClient} builds its TLS
context from the pinned CA as the sole root (system roots excluded), so
every daemon RPC --- CLI, \texttt{quarry mcp}, hooks, library --- verifies
against it.  The \texttt{quarry.toml} profile also records the CA under
\texttt{ca\_cert}, which \texttt{mcp-proxy} reads (as the sole root, TLS
1.3 minimum) when it transports MCP for other MCP clients; it is no longer
on quarry's own MCP path (DES-031 v2.2).

\subsection{Log Safety (CWE-532, CWE-117)}

Uvicorn's access log is disabled (\texttt{access\_log=False}),
eliminating query string leakage at the source.  The search handler
logs only result count, never the raw query.  (The WebSocket
session-key sanitization the prior design described was removed with the
in-daemon \texttt{/mcp} route --- there is no longer a WebSocket surface
to log.)

\subsection{Capture PII Redaction (DES-036)}

Captures --- session transcripts and WebFetch auto-captures --- are the
input to the planned private shadow-repo sync (quarry-ow3k), so
un-redacted content in a capture is a leak into a git-tracked (and soon
pushed) surface.  Every capture is scrubbed \textbf{at write time},
fail-closed, before any bytes reach disk.

\textbf{Scrubber (\modpath{scrub.py}).}  A \texttt{Scrubber} applies
redaction passes in a load-bearing order: \emph{secrets $\to$ paths
$\to$ emails $\to$ hostname $\to$ profanity}.  Email must precede
hostname so a hostname inside an email domain is subsumed by the
whole-email redaction rather than half-redacted (which would leak the
local part).  Each pass emits a marker no pass can re-match, so scrubbing
is idempotent --- backfill may safely re-run.

\begin{itemize}
  \item \textbf{Secrets} (\modpath{scrub\_rules.py}): GitHub PATs, AWS
    keys, Anthropic/OpenAI keys, bearer headers, JWTs, PEM/GPG private
    key blocks, \texttt{KEY=value} env secrets, Slack tokens.
  \item \textbf{Paths}: home directories (\texttt{/Users/<user>/},
    \texttt{/home/<user>/}) collapse to \texttt{\~{}}, retaining the
    useful deeper structure.
  \item \textbf{Emails}: any RFC-shaped address $\to$
    \texttt{[REDACTED:email]}.  A placeholder redacts every address
    (completeness beats attribution), not an ethos-handle mapping.
  \item \textbf{Hostname}: bounded to the local machine name
    (\texttt{socket.gethostname()} + \texttt{.local} + short leaf), not
    arbitrary dotted-token detection --- the latter has a catastrophic
    false-positive rate on identifiers like \texttt{github.com} or
    \texttt{quarry.db.facade}.
\end{itemize}

\textbf{Single write choke point (\modpath{capture.py}).}  Every local
\texttt{.md} capture producer --- \texttt{SessionTranscriptCapture}
(\modpath{session\_transcript.py}, serving the PreCompact, SessionEnd, and
SubagentStop hooks) and the backfill path (\modpath{backfill.py}) ---
routes through one \texttt{CaptureWriter}.  Scrubbing runs to completion before an atomic
temp-file rename, so a scrub failure writes no file and a partial or
half-redacted capture is never left behind.

\textbf{WebFetch DB-ingest.}  The WebFetch auto-capture path scrubs page
text via an opt-in \texttt{content\_scrubber} on the
\texttt{UrlIngest}/\texttt{InlineIngest} request
(\modpath{ingestion/web\_ingest.py}), which the daemon's
\texttt{CaptureIngestJob} and \texttt{ScrubbedIngestJob}
(\modpath{daemon/ingest\_jobs.py}) set for every hook-originated capture.
The field defaults to \texttt{None}: user-initiated \texttt{quarry ingest <url>},
sitemap crawls, and directory sync are \emph{not} scrubbed --- redaction
is a captures concern, not a general-ingestion one.  The URL's persisted
metadata form is redacted separately by \texttt{CaptureUrl}
(\modpath{capture\_url.py}), which strips userinfo, query, and fragment
(preserving IPv6 brackets) before the same text scrubber runs over the
bare \texttt{scheme://host/path}.

% ===================================================================
\section{Deployment}
% ===================================================================

\subsection{Single Machine}

\begin{lstlisting}
curl -fsSL https://raw.githubusercontent.com/.../install.sh | sh
\end{lstlisting}

Installs quarry, embedding model, mcp-proxy, Claude Code plugin, and
registers a daemon with TLS on localhost.

\subsection{Network Mode (Remote GPU)}

\begin{lstlisting}
export QUARRY_API_KEY=$(openssl rand -hex 32)
curl -fsSL https://raw.githubusercontent.com/.../install.sh | sh -s -- --network
\end{lstlisting}

Same as default, but binds daemon to \texttt{0.0.0.0:8420} instead of
localhost.  Requires \texttt{QUARRY\_API\_KEY}.  Installs quarry,
embedding model, generates TLS certificates, registers systemd service.
Detects NVIDIA GPUs and swaps \texttt{onnxruntime} for
\texttt{onnxruntime-gpu}.  Prints CA fingerprint for client TOFU.
Claude Code plugin installed if \texttt{claude} CLI is on PATH.

\subsection{Client (Remote Access)}

\begin{lstlisting}
curl -fsSL https://raw.githubusercontent.com/.../install.sh | sh
quarry login <server> --api-key <token>
\end{lstlisting}

No special flag needed.  Default install runs a local daemon on
localhost.  \texttt{quarry login} redirects queries to the remote
server via TOFU certificate pinning and writes the client profile
(\texttt{\~{}/.punt-labs/mcp-proxy/quarry.toml}: URL, pinned CA, bearer)
that \texttt{ClientConfig} resolves for every surface.

\subsection{Daemon (\texttt{quarryd})}

\begin{lstlisting}
QUARRY_API_KEY=<token> quarryd --port 8420 --host 0.0.0.0 --tls
\end{lstlisting}

Non-loopback hosts require an API key (the server refuses to start
without one).

Managed by \texttt{launchd} (macOS) or \texttt{systemd} (Linux) via
\texttt{quarry install}.  API key delivered via systemd
\texttt{EnvironmentFile} (Linux) or launchd
\texttt{EnvironmentVariables} (macOS) --- never in process arguments.

\subsection{Container (Fly.io)}

Multi-stage Dockerfile, three stages:

\begin{enumerate}
  \item \textbf{deps} (\texttt{python:3.13-slim}): installs
    \texttt{uv}, runs \texttt{uv sync -{}-frozen -{}-no-dev}.  Two-pass
    for Docker layer caching (deps first, then source).
  \item \textbf{model} (extends \texttt{deps}): downloads the
    embedding model at build time ($\sim$120\,MB, int8 quantized).
    Cached as a separate Docker layer.
  \item \textbf{runtime} (\texttt{python:3.13-slim}): copies
    \texttt{/app} from deps, copies model cache from model stage.
    \texttt{QUARRY\_ROOT=/data}, port 8080.
\end{enumerate}

LanceDB data lives on a Fly persistent volume at \texttt{/data}.
\texttt{QUARRY\_API\_KEY} is set via \texttt{fly secrets}.
Cold start is $\sim$5\,s with the baked-in model.

% ===================================================================
\section{Performance}
% ===================================================================

\subsection{Critical Path}

\begin{center}
\begin{tabular}{lrl}
\toprule
\textbf{Operation} & \textbf{Latency} & \textbf{Notes} \\
\midrule
Model load         & $\sim$250\,ms & One-time per process \\
Embedding (1 chunk) & $\sim$50\,ms  & Apple Silicon, CPU \\
LanceDB search     & $<$10\,ms     & Vector similarity + filter \\
End-to-end search  & $<$100\,ms    & After model is loaded \\
\bottomrule
\end{tabular}
\end{center}

\subsection{Fire-and-Forget Writes}

Side-effect operations return immediately with a \field{task\_id}; the
work runs on the daemon's serialized per-\texttt{RouteKey} queue
(DES-042).  This is necessary because:

\begin{enumerate}
  \item MCP tool calls block the LLM's response stream.  A 30-second
    ingest would block Claude from responding.  The 202 unblocks the
    caller immediately.
  \item Bounded, non-blocking admission (503 on a full queue) prevents
    the resource exhaustion that unbounded \texttt{asyncio.create\_task}
    ingests caused (the load-90 starvation, quarry-lnog).
  \item One in-flight \texttt{progressive\_insert} per LanceDB table
    (single writer per \texttt{RouteKey}) plus the embed
    \texttt{Semaphore} clamped to~1 eliminate the concurrent
    delete-then-add lost-update race, without snapshotting mutable
    per-request state.
\end{enumerate}

\subsection{Single Table Design}

All chunks live in one LanceDB table (\texttt{chunks}).  Document and
collection boundaries are columns, not separate tables.  This
simplifies cross-document search and avoids table proliferation.
Filtering by document/collection/page\_type/source\_format uses
LanceDB's built-in filter predicates on the vector search.

% ===================================================================
\section{Plugin and Hook System}
% ===================================================================

The Claude Code plugin (\texttt{plugin/hooks/hooks.json}, each script a
one-line \texttt{quarry-hook <event>} dispatched by
\modpath{\_hook\_entry.py}) registers hooks that fire on session events
and tool calls.  All hooks are fail-open: exceptions are logged and the
process exits~0.  Every hook is a thin daemon client --- no hook imports
the engine --- and every capture lands in the \texttt{<repo>-captures}
collection the daemon derives from the hook's \texttt{cwd}
(\texttt{default-captures} when the directory is unregistered), scrubbed
server-side before storage.  Each is independently toggleable under
\texttt{auto\_capture} in \texttt{.punt-labs/quarry/config.md}; only
\texttt{read} defaults to off.

\begin{center}
\begin{tabular}{llp{5cm}l}
\toprule
\textbf{Hook Event} & \textbf{Matcher} & \textbf{Effect} &
  \textbf{Blocks?} \\
\midrule
SessionStart  & * & Auto-register directory + background sync &
  No \\
PostToolUse   & quarry tools & Suppress tool output formatting &
  No \\
PostToolUse   & WebFetch & Auto-ingest fetched URL (scrubbed) &
  No \\
PostToolUse   & WebSearch & File a scrubbed digest of the results &
  No \\
PostToolUse   & Read & Opt-in: capture an out-of-tree prose file &
  No \\
PreCompact    & * & Capture transcript before compaction (scrubbed) &
  No \\
SessionEnd    & * & Capture the full transcript on every close &
  No \\
SubagentStop  & * & Archive subagent transcript; distill its report
  into \texttt{memory-<handle>} & No$^{\dagger}$ \\
\bottomrule
\end{tabular}
\end{center}

$^{\dagger}$Claude Code treats \texttt{SubagentStop} as a blocking hook:
a non-zero exit or a \texttt{decision} field keeps the subagent running.
The handler returns \texttt{\{\}} on every path, and its tests assert
that under adversarial payloads.

\subsection{Session Start}

\texttt{handle\_session\_start()} auto-registers the working directory
and launches a background sync via a detached subprocess.  Single-flight
is a kernel \texttt{flock} on \texttt{\~{}/.punt-labs/quarry/sync.pid}
(\texttt{SyncLock}, \modpath{sync\_lock.py}) held for the sync
subprocess's lifetime through fd inheritance, so a dead holder's lock is
released by the kernel with no PID parsing or reclaim race.  Returns
\texttt{additionalContext} immediately without waiting.

\subsection{Web Fetch Capture}

\texttt{handle\_post\_web\_fetch()} hands the already-fetched page to the
daemon's \texttt{/v1/capture} route, which extracts, scrubs, and stores
it in \texttt{<repo>-captures} (no second fetch; a JS-rendered page that
extracts to zero chunks is re-fetched server-side through the SSRF gate).
Before sending, \texttt{WebFetchLoopCloser}
(\modpath{web\_fetch\_loop\_closer.py}) checks \texttt{/v1/captures/lookup}
for a prior capture of the URL and, when one exists, nudges the agent to
\texttt{find} it instead of re-fetching.  The URL's persisted metadata form
is redacted by \texttt{CaptureUrl} (DES-036).

\subsection{Web Search and Read Capture}

\texttt{HookAgent.post\_web\_search()} files a scrubbed digest of a
WebSearch's results under the document name \texttt{search: <query>}.
\texttt{HookAgent.post\_read()} is opt-in (\texttt{read: false} by default,
because \texttt{Read} fires on every file the agent opens): a file passes
four admission checks in order --- outside every registered tree, not on
the secret-path denylist, an allow-listed extension, under the size cap
--- plus a client-side secret-pattern scan, before its text is sent.

\subsection{Pre-Compact and Session-End Capture}

\texttt{handle\_pre\_compact()} and \texttt{HookAgent.session\_end()} run
the same \texttt{SessionTranscriptCapture} pipeline
(\modpath{session\_transcript.py}): archive the raw JSONL under
\texttt{\~{}/.punt-labs/quarry/sessions/}, extract user/assistant text
(tool-use blocks skipped, capped at 500{,}000 characters), write the
scrubbed \texttt{.md} next to the project through the shared
\texttt{CaptureWriter}, and POST a \texttt{CaptureIngestRequest} the daemon
files as \texttt{session-<id8>} in \texttt{<repo>-captures}, attributed to
the repo's ethos pin.  PreCompact fires only on compaction; SessionEnd
fires on every close, so a short session that never compacts still yields
a durable capture.  A down daemon leaves the local archive for
\texttt{quarry backfill-sessions}.

\subsection{Subagent Stop Capture}

\texttt{HookAgent.subagent\_stop()} captures the subagent's own transcript
(\texttt{agent\_transcript\_path}, keyed by \texttt{agent\_id} --- the
parent's \texttt{transcript\_path}/\texttt{session\_id} are not its), then
--- only when \texttt{agent\_type} names a registered ethos identity and
the transcript had text --- files the subagent's final assistant turn as
an \texttt{observation} in \texttt{memory-<handle>} named
\texttt{subagent-<id8>-report}, through \texttt{/v1/remember} with the
same 5\,s cap as the first post (\modpath{subagent\_capture.py},
\modpath{subagent\_report.py}).  The transcript is parsed once for both
rows.  No model distills anything: the report is the summary the subagent
already wrote.  A non-identity \texttt{agent\_type}
(\texttt{general-purpose}) gets the raw capture only, never a row under
the leader's handle (DES-055).

% ===================================================================
\section{CLI Commands}
% ===================================================================

\begin{center}
\begin{longtable}{lp{8cm}}
\toprule
\textbf{Command} & \textbf{Purpose} \\
\midrule
\endfirsthead
\toprule
\textbf{Command} & \textbf{Purpose} \\
\midrule
\endhead
\texttt{ingest <url>}         & Ingest a URL (URL-only; local files and
  directories go through \texttt{register} + \texttt{sync}).  Flags:
  \texttt{-{}-agent-handle}, \texttt{-{}-memory-type}, \texttt{-{}-summary} \\
\texttt{remember}             & Ingest stdin content (\texttt{-{}-name};
  \texttt{-{}-agent-handle} routes to \texttt{memory-<handle>};
  \texttt{-{}-memory-type} from the closed vocabulary) \\
\texttt{learn <lesson>}       & Save a distilled lesson (\texttt{-{}-topic},
  \texttt{-{}-name}); the only writer of \texttt{memory\_type = lesson}
  (DES-053) \\
\texttt{find <query>}         & Semantic search \\
\texttt{show <document>}      & Document metadata or page text \\
\texttt{status}               & Database statistics \\
\texttt{use <name>}           & Set persistent default database \\
\texttt{delete <name>}        & Delete document or collection \\
\texttt{register <directory>} & Register for incremental sync \\
\texttt{deregister <collection>} & Remove directory registration \\
\texttt{sync}                 & Sync all registered directories \\
\texttt{enable [directory]}   & Enable quarry knowledge capture for a project directory \\
\texttt{disable [directory]}  & Disable quarry knowledge capture; remove registration and data \\
\texttt{optimize}             & Compact LanceDB table and rebuild indexes (CLI-only).
  Flags: \texttt{-{}-force} \\
\texttt{list documents}       & List indexed documents \\
\texttt{list collections}     & List collections with counts \\
\texttt{list databases}       & List named databases \\
\texttt{list registrations}   & List registered directories \\
\texttt{mcp}                  & Start MCP server (stdio; the engine runs
  only in \texttt{quarryd}) \\
\texttt{missions sync}        & File each frozen ethos mission round into
  \texttt{memory-<worker>} (\texttt{-{}-mission}, \texttt{-{}-dry-run},
  \texttt{-{}-force}; DES-055) \\
\texttt{captures init|push}   & Bootstrap / re-scrub and push the project's
  private capture shadow repo \\
\texttt{backfill-sessions}    & Backfill historical Claude Code session
  transcripts (returns 202) \\
\texttt{login <host>}         & TOFU login to remote server (cert pinning) \\
\texttt{logout}              & Disconnect from remote server \\
\texttt{remote list}         & Show remote config; \texttt{-{}-ping} validates connection \\
\texttt{install}             & Set up data directory, download model, generate TLS certs,
  register daemon, detect GPU, configure ethos ext \\
\texttt{doctor}              & Check environment health \\
\texttt{uninstall}           & Remove system daemon \\
\texttt{version}             & Print version \\
\bottomrule
\end{longtable}
\end{center}

Global flags: \texttt{-{}-json}, \texttt{-{}-verbose}/\texttt{-v},
\texttt{-{}-quiet}/\texttt{-q}, \texttt{-{}-db <name>}.

\textbf{Remote routing (DES-021):} When a remote server is configured
(via \texttt{quarry login}), all data commands should route to the
remote server over HTTPS with the pinned CA cert.  Only authentication
commands (\texttt{login}, \texttt{logout}, \texttt{remote}) and
administration commands (\texttt{install}, \texttt{uninstall},
\texttt{mcp}, \texttt{version}) are local-only.
All data commands --- \texttt{find}, \texttt{show}, \texttt{status},
\texttt{ingest}, \texttt{remember}, \texttt{learn}, \texttt{delete},
\texttt{sync}, every \texttt{list} subcommand, directory
\texttt{register}/\texttt{deregister}, and \texttt{missions sync} (which
reads the local mission files and posts each round to the configured
daemon) --- route remotely.  \texttt{use}
is local-only by design (the remote server is fixed to its own
database).

% ===================================================================
\section{Test Architecture}
% ===================================================================

The default run (\texttt{make test}) deselects the \texttt{slow} and
\texttt{embedding} tiers; \texttt{uv run pytest -{}-collect-only -q}
reports the current test count.  Three tiers plus integration; the files
named are representative, not exhaustive:

\begin{center}
\begin{tabular}{lp{5cm}p{5cm}}
\toprule
\textbf{Tier} & \textbf{Strategy} & \textbf{Files} \\
\midrule
Tier 1: Core logic & Real data, no mocks & \texttt{test\_database},
  \texttt{test\_chunker}, \texttt{test\_embeddings},
  \texttt{test\_code\_processor}, \texttt{test\_html\_processor},
  \texttt{test\_text\_processor}, \texttt{test\_collections},
  \texttt{test\_config}, \texttt{test\_formatting},
  \texttt{test\_models}, \texttt{test\_retrieval\_fusion},
  \texttt{test\_retrieval\_hybrid}, \texttt{test\_memory\_types},
  \texttt{test\_scrub}, \texttt{test\_format\_strategies},
  \texttt{test\_registry} \\
Tier 2: Orchestration & Mocked I/O & \texttt{test\_pipeline},
  \texttt{test\_pipeline\_images},
  \texttt{test\_ocr\_local}, \texttt{test\_url\_ingestion},
  \texttt{test\_sitemap}, \texttt{test\_backends},
  \texttt{test\_sync}, \texttt{test\_sync\_concurrency},
  \texttt{test\_hooks}, \texttt{test\_hooks\_agent},
  \texttt{test\_subagent\_capture}, \texttt{test\_mission\_memory},
  \texttt{test\_mission\_store}, \texttt{test\_doctor} \\
Tier 3: Surface wiring & Fully mocked / hermetic in-process ASGI
  (\texttt{httpx.ASGITransport}), $<$2\,s &
  \texttt{test\_cli}, \texttt{test\_cli\_missions},
  \texttt{test\_mcp\_server}, \texttt{test\_mcp\_missions},
  \texttt{test\_http\_server} (the \modpath{api/} +
  \modpath{daemon/routes/} contract), \texttt{test\_api},
  \texttt{test\_client}, \texttt{test\_client\_config},
  \texttt{test\_retrieval\_equivalence},
  \texttt{test\_daemon\_launcher}, \texttt{test\_watch\_loop},
  \texttt{test\_watch\_reconcile}, \texttt{test\_fs\_watchdog},
  \texttt{test\_proxy}, \texttt{test\_service} \\
Invariants & Long-lived-process leak guards; suite hermeticity
  (DES-047) & \texttt{test\_resource\_invariants},
  \texttt{test\_hermeticity} \\
Integration & Real ONNX + real LanceDB; real loopback TLS &
  \texttt{test\_integration}, \texttt{test\_tls\_smoke}
  (\texttt{@pytest.mark.slow}; the model loads only under
  \texttt{embedding}) \\
\bottomrule
\end{tabular}
\end{center}

Quality gates (\texttt{make check}): \texttt{ruff} (lint + format),
\texttt{mypy} (strict), \texttt{pyright} (strict), \texttt{pytest} (all
markers except slow/embedding), the OO, coupling, and suppression
merge-base ratchets (DES-040), the import-linter client/engine contract,
and the OpenAPI drift check.  Zero violations required.

% ===================================================================
\section{Formal Model}
% ===================================================================

The plugin state machine --- how quarry hooks into the Claude Code
session lifecycle --- is formally specified in
\texttt{docs/claude-code-quarry.tex} using Z notation (ISO~13568).

The specification covers:

\begin{itemize}
  \item Session start with/without auto-sync
  \item Web fetch auto-capture into the \texttt{<repo>-captures} collection
  \item Pre-compaction transcript capture into \texttt{<repo>-captures}
  \item Database switching via \texttt{use} tool
  \item Knowledge hint injection at tool boundaries (specified, not yet
    built)
  \item Compaction observability: a capture flag records whether
    capture preceded compaction; the hook system enforces
    capture-before-compaction on a best-effort basis
  \item Clean session boundaries: all quarry state cleared on
    \texttt{EndSession} --- now wired to the \texttt{SessionEnd} hook,
    which also captures the full transcript
  \item The WebSearch, Read, and SubagentStop captures ride base-layer
    operations (\texttt{CompleteToolSuccess}, \texttt{StopSubagent}); the
    spec's implementation-validation section records them as wired
    rather than modelling them as quarry-layer schemas
\end{itemize}

The specification can be type-checked with \texttt{fuzz -t
claude-code-quarry.tex}.

% ===================================================================
\appendix
% ===================================================================

\section{Operation Concurrency Model}
\label{app:concurrency}

Every quarry operation is exposed on up to three surfaces: the HTTP
API (\texttt{/v1}, daemon), the CLI, and the MCP server (\texttt{quarry
mcp} stdio).  The concurrency model differs by surface because each has
different constraints: HTTP connections time out, MCP tool calls block
the LLM response stream, and CLI commands run in a user's terminal.
Under daemon-first every CLI and MCP call is a \texttt{QuarryClient}
request to the daemon, so the ``local'' and ``remote'' CLI columns below
differ only in transport target, not in code path.

Table~\ref{tab:concurrency} classifies each operation.  The
\textbf{Client model} column describes what the caller experiences:

\begin{description}[style=nextline,leftmargin=2.5cm]
  \item[sync] The caller blocks until the operation completes and
    receives the result inline.
  \item[fire-and-forget] The caller receives an acknowledgement
    immediately (HTTP~202 or a short string).  Work runs in the
    background; the caller polls or ignores the outcome.
  \item[---] The operation is not available on that surface.
\end{description}

\subsection{Concurrency Matrix}

\begin{longtable}{l l l l l}
\caption{Operation concurrency by surface.
  \emph{Long-running} = involves network~I/O, embedding, or bulk
  LanceDB writes.}\label{tab:concurrency} \\
\toprule
\textbf{Operation} & \textbf{HTTP API} & \textbf{CLI (local)}
  & \textbf{CLI (remote)} & \textbf{MCP} \\
\midrule
\endfirsthead
\toprule
\textbf{Operation} & \textbf{HTTP API} & \textbf{CLI (local)}
  & \textbf{CLI (remote)} & \textbf{MCP} \\
\midrule
\endhead
\bottomrule
\endfoot
%
% --- Read operations (fast, sync everywhere) ---
\multicolumn{5}{l}{\emph{Read operations}} \\
\midrule
find (search)       & sync   & sync & sync & sync \\
show                & sync   & sync & sync & sync \\
status              & sync   & sync & sync & sync \\
list documents      & sync   & sync & sync & sync \\
list collections    & sync   & sync & sync & sync \\
list registrations  & sync   & sync & sync & sync \\
list databases      & sync   & sync & sync & sync \\
health              & sync   & ---  & ---  & ---  \\
\midrule
%
% --- Mutating operations (all async on HTTP) ---
\multicolumn{5}{l}{\emph{Mutating operations}} \\
\midrule
register            & fire-and-forget & fire-and-forget & fire-and-forget & fire-and-forget \\
deregister          & fire-and-forget & fire-and-forget & fire-and-forget & fire-and-forget \\
delete (document)   & fire-and-forget & fire-and-forget & fire-and-forget & fire-and-forget \\
delete (collection) & fire-and-forget & fire-and-forget & fire-and-forget & fire-and-forget \\
remember            & fire-and-forget & fire-and-forget & fire-and-forget & fire-and-forget \\
learn               & fire-and-forget & fire-and-forget & fire-and-forget & fire-and-forget \\
ingest (URL)        & fire-and-forget & fire-and-forget & fire-and-forget
                    & fire-and-forget \\
capture (hooks)     & fire-and-forget & ---  & ---  & --- \\
sync                & fire-and-forget & fire-and-forget & fire-and-forget
                    & fire-and-forget \\
missions sync       & --- (N $\times$ remember) & sync & sync & sync \\
use (switch db)     & sync   & sync & sync & sync \\
optimize            & fire-and-forget & fire-and-forget & fire-and-forget & --- \\
backfill-sessions   & fire-and-forget & fire-and-forget & fire-and-forget & --- \\
\end{longtable}

\subsection{Design Rationale}

\textbf{Read operations} are synchronous on all surfaces.  Hybrid
search (embedding + vector similarity + BM25) completes in
$<$100\,ms after model load; the overhead of task tracking would
exceed the operation itself.

\textbf{Mutating operations} are all fire-and-forget on HTTP, and
therefore on every client of it: the CLI prints the daemon's
\field{task\_id} and exits~0.  Every mutating endpoint returns
202~Accepted with a \field{task\_id}; the work runs on the daemon's
serialized per-\texttt{RouteKey} ingest queue (DES-042).  This
eliminates the 15-second CLI timeout as a failure mode for any write
operation.  Local files never travel over the wire: \texttt{ingest} is
URL-only on every surface, and a directory is indexed by
\texttt{register} + the daemon's own watch loop.  \texttt{missions sync}
is the one client-side composite: it reads the mission files locally and
issues one \texttt{remember} per frozen round, so it has no HTTP row of
its own.

\begin{itemize}
  \item All mutating endpoints use the single unified
    \texttt{GET /v1/tasks/\{task\_id\}} for status polling; the old
    \texttt{/sync/\{id\}} and \texttt{/ingest/\{id\}} alias routes were
    removed in the DES-031 v2.2 contract rewrite.
  \item \textbf{Sync} no longer returns 409: it enqueues transparently
    (202 + poll), because the daemon's serialized per-\texttt{RouteKey}
    queue (DES-042) already prevents concurrent writers to a table.  The
    409 is retained only for \textbf{optimize} and
    \textbf{backfill-sessions}, whose whole-table compaction/rebuild
    must not overlap itself.
  \item Completed and failed tasks are garbage-collected after a
    1-hour TTL on the next task creation.
  \item \textbf{Optimize} and \textbf{backfill-sessions} are exposed on
    HTTP and the CLI but \emph{not} as MCP tools --- they are
    whole-database maintenance, run with explicit intent rather than
    from an agent's tool loop.
\end{itemize}

\subsection{Task Lifecycle}

All background tasks follow the same lifecycle:

\begin{enumerate}
  \item \texttt{POST} validates input (including the closed
    \texttt{memory\_type} vocabulary on the write routes), creates a
    \texttt{TaskState} record in the daemon's \texttt{TaskRegistry}
    (\modpath{daemon/tasks.py}), hands the job to the daemon's background
    machinery (the per-\texttt{RouteKey} ingest queue for writes), and
    returns HTTP~202 with \field{task\_id} and
    \field{status}~=~\texttt{accepted}.
  \item The background task updates \field{status} to
    \texttt{completed} (with \field{results}) or \texttt{failed}
    (with \field{error}) on completion.
  \item A \texttt{finally} guard marks any task that exits without
    setting a terminal status as \texttt{failed}.
  \item \texttt{CancelledError} is caught before \texttt{Exception}
    to ensure clean shutdown on server stop.
  \item The CLI prints the \field{task\_id} and exits~0 (fire-and-forget).
\end{enumerate}

\end{document}
